%% ============================================================================
%%  main.tex
%%  قالب رسمی پایان‌نامه‌ها و رساله‌های دانشگاه خوارزمی
%%  Kharazmi University Official Persian LaTeX Thesis Template
%% ============================================================================
%%
%%  راهنمای سریع کامپایل:
%%  این قالب باید با موتور XeLaTeX همراه با BibTeX و MakeIndex کامپایل شود.
%%  ترتیب دستورات کامپایل در خط فرمان:
%%      xelatex -interaction=nonstopmode main.tex
%%      bibtex main
%%      makeindex main
%%      xelatex -interaction=nonstopmode main.tex
%%      xelatex -interaction=nonstopmode main.tex
%%
%%  گزینه‌های کلاس kharazmi-thesis:
%%      msc               : پایان‌نامه کارشناسی ارشد (پیش‌فرض)
%%      phd               : رساله دکتری تخصصی
%%      bsc               : پروژه کارشناسی
%%      harfifrontmatter  : شماره‌گذاری حروفی صفحات آغازین (اختیاری)
%%
%% ============================================================================

\documentclass[msc]{kharazmi-thesis}

% بارگذاری بسته تنظیمات ظاهری، فونت‌ها و فراداده‌های دانشگاه خوارزمی
% گزینه‌های فونت: bzar (پیش‌فرض مصوب)، bnazanin، vazirmatn
\usepackage{kharazmi}

% فعال‌سازی سامانه نمایه موضوعی و الفبایی
\makeindex

% -----------------------------------------------------------------------------
% ۱. تنظیمات مشخصات و فراداده‌های پایان‌نامه (Metadata)
% -----------------------------------------------------------------------------
\university{دانشگاه خوارزمی}
\universityen{Kharazmi University}

\ministry{وزارت علوم، تحقیقات و فناوری}
\ministryen{Ministry of Science, Research and Technology}

\faculty{دانشکده مهندسی برق و کامپیوتر}
\facultyen{Faculty of Electrical and Computer Engineering}

\department{گروه مهندسی کامپیوتر}
\departmenten{Department of Computer Engineering}

% عنوان پایان‌نامه به فارسی و انگلیسی
\thesistitle{طراحی و پیاده‌سازی قالب جامع پایان‌نامه دانشگاه خوارزمی در سامانه لاتک}
\thesistitleen{Design and Implementation of Kharazmi University Comprehensive LaTeX Thesis Template}

% مشخصات نویسنده
\authorname{آیریک پژوهشگر}
\authornameen{Ayrik Researcher}
\studentid{1402123456}

% مقطع و رشته تحصیلی
\field{مهندسی کامپیوتر}
\fielden{Computer Engineering}
\major{هوش مصنوعی}
\majoren{Artificial Intelligence}

% استادان راهنما و مشاور
\supervisor[دانشیار]{دکتر احمد رضایی}
\supervisoren{Dr. Ahmad Rezaei}

\advisor[استادیار]{دکتر سارا محمدی}
\advisoren{Dr. Sara Mohammadi}

% اعضای هیئت داوران جلسه دفاع
\internaljudge[استاد]{دکتر علی علوی}
\internaljudgeen[Professor]{Dr. Ali Alavi}
\externaljudge[دانشیار]{دکتر رضا اکبری}{دانشگاه تهران}
\externaljudgeen[Associate Professor]{Dr. Reza Akbari}{University of Tehran}

% تاریخ دفاع و فارغ‌التحصیلی
\thesisdate{شهریور ۱۴۰۳}
\thesisdateen{September 2024}

% کلیدواژه‌ها به زبان فارسی و انگلیسی
\keywords{قالب پایان‌نامه، دانشگاه خوارزمی، زی‌پرشین، حروف‌چینی علمی، مراجع دوزبانه.}
\keywordsen{Thesis Template, Kharazmi University, XePersian, Academic Typesetting, Bilingual Bibliography.}

% =============================================================================
% شروع بدنه اصلی سند
% =============================================================================
\begin{document}

% -----------------------------------------------------------------------------
% بخش اول: صفحات آغازین پیش از متن (Front Matter)
% ترتیب صفحات بر اساس جدول ۲-۱ شیوه‌نامه نگارش دانشگاه خوارزمی:
% جلد -> بسم‌الله (صفحه ۱) -> تعهدنامه -> تصویب‌نامه -> تقدیم -> سپاسگزاری -> چکیده
% -----------------------------------------------------------------------------
%% ============================================================================
%%  front/cover_fa.tex
%%  صفحه جلد فارسی پایان‌نامه (برگ نخست)
%%  طراحی شده برای قالب رسمی دانشگاه خوارزمی
%% ============================================================================
%%
%%  توضیح:
%%  این فایل صفحه جلد فارسی پایان‌نامه را تولید می‌نماید.
%%  اطلاعات عنوان، نام دانشجو، استادان راهنما و مشاور، مقطع و تاریخ
%%  به‌طور خودکار از متغیرهای تعریف‌شده در فایل main.tex فراخوانی می‌شوند.
%%  در صورت تمایل به تغییر ظاهر یا شخصی‌سازی، می‌توانید مستقیماً این فایل را ویرایش نمایید.
%% ============================================================================

\makefaCover
        % ۱. صفحه جلد فارسی

\frontmatter                      % شروع شماره‌گذاری پیوسته از ۱ از صفحه بعد

%% ============================================================================
%%  front/besmellah.tex
%%  صفحه بسم‌الله الرحمن الرحیم (برگ دوم)
%%  طراحی شده برای قالب رسمی دانشگاه خوارزمی
%% ============================================================================
%%
%%  توضیح:
%%  این فایل صفحه بسم‌الله پایان‌نامه را تولید می‌کند.
%%  در صورتی که فایل تصویری images/besmellah.png وجود داشته باشد، آن تصویر را درج می‌نماید؛
%%  در غیر این صورت از خط زیبای نستعلیق با عبارت «بسم الله الرحمن الرحیم» استفاده می‌شود.
%% ============================================================================

\makeBesm
       % ۲. صفحه بسم‌الله الرحمن الرحیم (صفحه ۱)
%% ============================================================================
%%  front/declaration.tex
%%  تعهدنامه اصالت و مالکیت اثر
%%  طراحی شده برای قالب رسمی دانشگاه خوارزمی
%% ============================================================================
%%
%%  توضیح:
%%  این برگه بر اساس ضوابط آموزشی و پژوهشی دانشگاه خوارزمی تنظیم شده و بیانگر
%%  اصالت علمی پژوهش، رعایت حقوق معنوی دیگران و مالکیت مادی دانشگاه نسبت به اثر است.
%%  اطلاعات نام دانشجو، رشته، گرایش و مقطع به‌صورت خودکار در متن درج می‌شوند.
%% ============================================================================

\makeDeclaration
     % ۳. تعهدنامه اصالت و مالکیت اثر
%% ============================================================================
%%  front/approval.tex
%%  صورت‌جلسه دفاع و برگ تأیید هیئت داوران
%%  طراحی شده برای قالب رسمی دانشگاه خوارزمی
%% ============================================================================
%%
%%  توضیح:
%%  این فایل برگه تصویب و صورت‌جلسه دفاع پایان‌نامه را بر اساس مشخصات
%%  استادان راهنما، مشاور و داوران داخلی و خارجی تولید می‌نماید.
%%  اطلاعات اعضای هیئت داوران در فایل اصلی main.tex قابل تنظیم است.
%%  پس از جلسه دفاع و دریافت امضاهای نهایی، می‌توانید صفحه اسکن‌شده با کیفیت
%%  را جایگزین این برگه نمایید یا از همین خروجی استاندارد استفاده کنید.
%% ============================================================================

\makeApproval
        % ۴. صورت‌جلسه دفاع و تصویب‌نامه هیئت داوران
%% ============================================================================
%%  front/dedication.tex
%%  صفحه تقدیم پایان‌نامه
%%  طراحی شده برای قالب رسمی دانشگاه خوارزمی
%% ============================================================================
%%
%%  توضیح:
%%  این صفحه برای تقدیم اثر به اشخاص و مراجع مورد نظر دانشجو اختصاص دارد.
%%  شما می‌توانید متن زیر را به دلخواه خود ویرایش و شخصی‌سازی نمایید.
%% ============================================================================

\makeDedication{%
  تقدیم به پدر و مادر مهربان و فداکارم،\\
  که مهر بی‌پایان، دعای خیر و فداکاری‌های بی‌دریغشان\\
  همواره روشنایی‌بخش مسیر زندگی و تحصیلم بوده است؛\\[1.5em]
  به همسر همراه و صبورم که با شکیبایی مرا در این مسیر یاری نمود؛\\[1.5em]
  و به تمامی اساتید و آموزگاران فرزانه‌ای که\\
  به من درس دانش، اخلاق، و انسانیت آموختند.
}
      % ۵. صفحه تقدیم
%% ============================================================================
%%  front/acknowledgments.tex
%%  صفحه سپاسگزاری پایان‌نامه
%%  طراحی شده برای قالب رسمی دانشگاه خوارزمی
%% ============================================================================
%%
%%  توضیح:
%%  این صفحه برای قدردانی و سپاسگزاری از استادان راهنما، مشاور، داوران،
%%  نهادهای حامی پژوهش و خانواده در نظر گرفته شده است.
%%  می‌توانید متن زیر را با توجه به شرایط و همراهان پژوهش خود ویرایش فرمایید.
%% ============================================================================

\makeAck{%
  سپاس بی‌کران پروردگار یکتا را سزاست که به من توفیق، دانایی و انگیزه پیمودن این مسیر علمی را عنایت فرمود.\\[1em]
  بر خود لازم و شایسته می‌دانم از استاد ارجمند و راهنمای فرزانه‌ام، جناب آقای/سرکار خانم دکتر ...، به پاس راهنمایی‌های دلسوزانه و بی‌دریغ، دقت علمی موشکافانه، و صبر و گشاده‌رویی در تمامی مراحل انجام این پژوهش، صمیمانه‌ترین سپاس‌ها و احترامات قلبی خود را تقدیم نمایم.\\[1em]
  همچنین از استاد(ان) مشاور محترم و اعضای گرامی هیئت داوران که با مطالعه ژرف و دقیق این رساله/پایان‌نامه و ارائه نکات و انتقادات سازنده، موجبات غنای علمی این اثر را فراهم آوردند، کمال امتنان و تشکر را دارم.\\[1em]
  از مسئولان، کارشناسان و کادر محترم دانشکده مهندسی برق و کامپیوتر دانشگاه خوارزمی که امکانات پژوهشی و آزمایشگاهی لازم را در اختیار اینجانب قرار دادند، قدردانی می‌نمایم.\\[1em]
  در نهایت، از پدر و مادر عزیزم و اعضای محترم خانواده‌ام که با مهر بی‌کران، شکیبایی و دعای خیر خود همواره بزرگ‌ترین حامی و مشوق من در لحظات دشوار تحقیق بودند، صمیمانه سپاسگزاری می‌نمایم.
}
 % ۶. صفحه سپاسگزاری
%% ============================================================================
%%  front/abstract_fa.tex
%%  چکیده پایان‌نامه به زبان فارسی و کلیدواژه‌ها
%%  طراحی شده برای قالب رسمی دانشگاه خوارزمی
%% ============================================================================
%%
%%  توضیح:
%%  بر اساس شیوه‌نامه نگارش دانشگاه خوارزمی:
%%  - متن چکیده باید شامل بیان هدف، روش پژوهش، و یافته‌های کلیدی باشد.
%%  - تعداد کلمات چکیده معمولاً بین ۲۰۰ تا ۳۰۰ کلمه است.
%%  - در پایان چکیده، کمینه ۳ و بیشینه ۵ کلیدواژه قید می‌گردد.
%% ============================================================================

\abstractfa{%
امروزه تدوین و صفحه‌آرایی مستندات و رساله‌های علمی با رعایت دقیق استانداردهای نگارشی دانشگاهی نقشی اساسی در ارتقای کیفیت آثار پژوهشی ایفا می‌کند. در این پژوهش، یک سامانه جامع و استاندارد برای حروف‌چینی پایان‌نامه‌ها و رساله‌های دانشگاه خوارزمی بر پایه سامانه \lr{\XeLaTeX} و بسته زی‌پرشین طراحی و پیاده‌سازی شده است. در این ساختار، کلیه ضوابط مصوب دانشگاه شامل فواصل حاشیه‌ها (بالا ۳/۵، پایین ۲/۵، راست ۳/۰ و چپ ۳/۰ سانتی‌متر)، قلم‌های رسمی بی زر و بی تیتر، اندازه و فواصل دقیق سرعنوان‌ها، شماره‌گذاری فارسی تمامی عناصر (بخش‌ها، جدول‌ها، تصویرها، نمودارها و روابط ریاضی)، سرصفحه‌های متقارن، پانویس‌های استاندارد هر صفحه، و تفکیک منظم اجزای پیش از متن و پس از متن به‌طور دقیق رعایت شده است. علاوه‌بر این، سازوکار مدیریت مراجع دوزبانه (فارسی و انگلیسی) با استفاده از شیوه‌نامه‌های استاندارد بیب‌تک پیاده‌سازی شده تا ارجاع‌دهی به مقالات و کتاب‌ها بدون تداخل در جهت متن انجام پذیرد. نتایج ارزیابی‌های فنی نشان می‌دهد که قالب حاصل ضمن پایداری کامل در فرآیند کامپایل و فقدان هرگونه نویسه گم‌شده، کاربری آسانی برای دانشجویان تحصیلات تکمیلی به همراه دارد.
}

\keywords{قالب پایان‌نامه لاتک، زی‌پرشین، دانشگاه خوارزمی، مراجع دوزبانه، استانداردهای نگارشی.}

\makeAbstractFa
     % ۷. چکیده فارسی و کلیدواژه‌ها

% -----------------------------------------------------------------------------
% فهرست‌های خودکار و جداول نشانه‌ها (Lists)
% -----------------------------------------------------------------------------
\tableofcontents                  % فهرست مطالب
\listoffigures                    % فهرست شکل‌ها
\listoftables                     % فهرست جدول‌ها
\listofcharts                     % فهرست نمودارها

% فهرست نشانه‌ها و نمادهای ریاضی (اختیاری)
\listofsymbols{
\symb{\alpha}{نرخ یادگیری در الگوریتم‌های بهینه‌سازی}
\symb{\beta}{ضریب تکانه گرادیان در میانگین متحرک}
\symb{\nabla f(x)}{بردار گرادیان تابع هدف در نقطه $x$}
\symb{\mathcal{L}(\theta)}{تابع زیان تجربی روی نمونه‌های آموزشی}
\symb{\mathbf{Q}, \mathbf{K}, \mathbf{V}}{ماتریس‌های پرسش، کلید و مقدار در لایه توجه خودکار}
}

% فهرست کوته‌نوشت‌ها و کلمات اختصاری (اختیاری)
\listofabbreviations{
\persianToEnglish{هوش مصنوعی}{AI: Artificial Intelligence}
\persianToEnglish{پردازش زبان طبیعی}{NLP: Natural Language Processing}
\persianToEnglish{شبکه عصبی پیچشی}{CNN: Convolutional Neural Network}
\persianToEnglish{شبکه عصبی عمیق}{DNN: Deep Neural Network}
\persianToEnglish{مدل زبان بزرگ}{LLM: Large Language Model}
}

% -----------------------------------------------------------------------------
% بخش دوم: فصول پنج‌گانه متن اصلی (Main Matter)
% -----------------------------------------------------------------------------
\mainmatter
%% ============================================================================
%%  chapters/chapter1.tex
%%  فصل ۱: کلیات پژوهش
%%  قالب رسمی پایان‌نامه و رساله دانشگاه خوارزمی
%% ============================================================================

\chapter{کلیات پژوهش}
\label{chap:intro}

\section{مقدمه}
\label{sec:ch1_intro}

توسعه روزافزون دانش و فناوری در جهان امروز، ساماندهی، مستندسازی و نشر دقیق یافته‌های علمی را به یکی از ارکان بنیادین پیشرفت جوامع علمی تبدیل کرده است. در این میان، پایان‌نامه‌ها و رساله‌های دانشگاهی، عصاره فعالیت‌های پژوهشی دانشجویان تحصیلات تکمیلی و بازتاب‌دهنده توان علمی دانشگاه‌ها و مراکز تحقیقاتی به شمار می‌روند. ارائه یک اثر علمی مستلزم رعایت دو مؤلفه اساسی است: محتوای غنی علمی و قالب استاندارد نگارشی. هرچند ارزش ذاتی یک رساله در گرو نوآوری‌های نظری، دقت در روش‌شناسی و صحت نتایج تجربی آن است، اما شیوه نمایش، صفحه‌آرایی منظم، یکدستی علائم، و رعایت دقیق دستورالعمل‌های نگارشی دانشگاه نقش تعیین‌کننده‌ای در انتقال صحیح پیام علمی و خوانایی اثر ایفا می‌کند.

در دهه‌های اخیر، سامانه‌های حروف‌چینی دیجیتال تحولی ژرف در تولید مستندات علمی پدید آورده‌اند. سامانه حروف‌چینی تک\index{تک@\lr{TeX}}\LTRfootnote{\lr{TeX}} که توسط دانلد کنوت ابداع گردید و به دنبال آن بسته توسعه‌یافته لاتک\index{لاتک@\lr{LaTeX}}\LTRfootnote{\lr{LaTeX}}، به عنوان ابزارهای معیار در جامعه علمی شناخته شده‌اند. قابلیت‌های منحصر‌به‌فرد این سامانه‌ها در حروف‌چینی دیجیتال\index{حروف‌چینی دیجیتال} فرمول‌های پیچیده ریاضی، ارجاع‌دهی خودکار، مدیریت یکپارچه فهرست‌ها و حفظ ثبات ساختاری متن، موجب تمایز آشکار آن‌ها از واژه‌پردازهای سنتی گردیده است. با این حال، استفاده از این سامانه‌ها در زبان‌هایی با جهت نگارش راست‌به‌چپ\LTRfootnote{\lr{Right-to-Left (RTL)}} نظیر زبان فارسی، همواره با چالش‌های فنی و ساختاری فراوانی روبه‌رو بوده است. ظهور موتور زی‌تک\LTRfootnote{\lr{XeTeX}} و توسعه بسته بومی زی‌پرشین\index{زی‌پرشین@\lr{XePersian}}\LTRfootnote{\lr{XePersian}} زمینه مناسبی را برای حل ریشه‌ای این مسائل و انطباق کامل با ظرافت‌های تایپوگرافی خط فارسی فراهم ساخت.

دانشگاه خوارزمی به عنوان نخستین مؤسسه آموزش عالی مدرن کشور، با تدوین شیوه‌نامه جامع نگارش پایان‌نامه‌ها و رساله‌ها\index{شیوه‌نامه دانشگاه خوارزمی}، استانداردهای دقیقی را برای ساختار، اندازه‌ها، قلم‌ها، حاشیه‌ها و صفحه‌آرایی آثار علمی تعیین کرده است. هدف اصلی این پژوهش، طراحی و پیاده‌سازی یک قالب پایدار، جامع و کاملاً منطبق بر شیوه‌نامه مصوب دانشگاه خوارزمی است تا پژوهشگران بتوانند بدون دغدغه‌های مربوط به تنظیمات ظاهری، بر کیفیت محتوای علمی خود تمرکز نمایند.

\section{بیان مسئله}
\label{sec:ch1_problem}

نگارش متون علمی دانشگاهی به زبان فارسی در بستر سامانه‌های دیجیتال همواره با پیچیدگی‌های چندگانه‌ای روبه‌رو بوده است. از یک سو، واژه‌پردازهای تجاری نظیر مایکروسافت ورد\LTRfootnote{\lr{Microsoft Word}} اگرچه دارای رابط کاربری بصری هستند، اما در مدیریت مستندات حجیم، مدیریت خودکار مراجع دوزبانه، شماره‌گذاری پایدار روابط ریاضی و جدول‌ها، و کنترل دقیق فواصل ساختاری دچار ناپایداری‌های مکرر می‌شوند. به‌هم‌ریختگی جایگاه تصویرها، تغییر ناخواسته قلم‌ها، و عدم تطابق شماره‌گذاری بخش‌ها در اسناد بزرگ از معضلات شناخته‌شده این ابزارهاست.

از سوی دیگر، حروف‌چینی متون فارسی در سامانه لاتک بدون وجود یک قالب دانشگاهی بومی‌سازی‌شده، با دشواری‌های فنی متعددی مواجه است:
\begin{enumerate}
  \item \textbf{تداخل جهت متن دوزبانه:} متون علمی مهندسی و علوم پایه همواره آکنده از اصطلاحات فنی لاتین، اختصارات انگلیسی، و روابط ریاضی در دل پاراگراف‌های فارسی هستند. مدیریت صحیح جهت جریان متن\LTRfootnote{\lr{Bidirectional Text Flow}} به‌گونه‌ای که علائم نگارشی نظیر نقطه، خط‌فاصله، پرانتز و پانویس‌ها در جایگاه درست خود قرار گیرند، نیازمند تنظیم دقیق قواعد دوسویه است.
  \item \textbf{شیوه‌نامه دقیق دانشگاه خوارزمی:} شیوه‌نامه مصوب دانشگاه دارای ضوابط بسیار دقیق در خصوص فواصل حاشیه‌ها (بالا ۳/۵، پایین ۲/۵، راست ۳/۰ و چپ ۳/۰ سانتی‌متر)، فاصله سرصفحه تا لبه کاغذ (۲/۲۵ سانتی‌متر)، اندازه قلم متن اصلی (بی زر ۱۳ با فاصله خطوط ۱/۱۵)، اندازه و فاصله سرعنوان‌ها، و سبک سرصفحه‌ها با ساختار «عنوان \lr{|} شماره صفحه» است که پیاده‌سازی دستی آن توسط دانشجویان بسیار وقت‌گیر و همراه با خطاست.
  \item \textbf{یکپارچگی شماره‌گذاری فارسی:} تمامی عناصر سند شامل شماره صفحه‌ها، شماره بخش‌ها، جدول‌ها، شکل‌ها، نمودارها، روابط ریاضی، الگوریتم‌ها و پانویس‌ها باید از ارقام فارسی تبعیت نمایند.
  \item \textbf{مدیریت مراجع دوزبانه:} درج هم‌زمان مراجع فارسی و لاتین بدون تداخل در سبک کتاب‌شناسی، مرتب‌سازی صحیح و حفظ الگوهای ارجاع‌دهی درون‌متنی، نیازمند سبک‌های اختصاصی بیب‌تک\LTRfootnote{\lr{BibTeX}} است.
\end{enumerate}

بنابراین، عدم وجود یک قالب رسمی، ماژولار و آزموده‌شده برای دانشگاه خوارزمی سبب اتلاف زمان قابل‌توجهی از دانشجویان تحصیلات تکمیلی در مراحل پایانی فراغت از تحصیل و بروز ناهمگونی در انتشارات دانشگاهی می‌گردد.

\section{اهداف پژوهش}
\label{sec:ch1_objectives}

هدف بنیادین این پژوهش، طراحی، پیاده‌سازی و اعتبارسنجی یک بسته و کلاس جامع لاتک برای پایان‌نامه‌ها و رساله‌های دانشگاه خوارزمی است. در راستای تحقق این هدف کلان، اهداف فرعی زیر تعریف شده‌اند:
\begin{itemize}
  \item طراحی کلاس اختصاصی \lr{\texttt{kharazmi-thesis.cls}} با بهره‌گیری از کلاس پایه \lr{\texttt{report}} برای اعمال مستقیم هندسه صفحه، فواصل، سرصفحه‌ها و قلم‌ها.
  \item پیاده‌سازی بسته تکمیلی \lr{\texttt{kharazmi.sty}} جهت تعریف متغیرهای فراداده، صفحات پیش از متن (جلد فارسی، تصویب‌نامه، تعهدنامه، تقدیم، سپاسگزاری، چکیده) و صفحات پس از متن (واژه‌نامه‌ها، کارنامک، فهرست مقالات، جلد انگلیسی).
  \item ایجاد سازوکار مدیریت مراجع دوزبانه با انطباق کامل با موتور زی‌پرشین.
  \item ارائه ۵ فصل نمونه واقعی با پوشش کامل فرمول‌های ریاضی، جداول تخصصی، تصویرها، نمودارها، الگوریتم‌ها و قطعه کدها.
  \item تدوین راهنمای جامع کاربری و مستندات گام‌به‌گام نصب و کامپایل برای استفاده دانشجویان.
\end{itemize}

\section{پرسش‌ها و فرضیه‌های پژوهش}
\label{sec:ch1_questions}

در راستای هدایت دقیق فرآیند طراحی و ارزیابی فنی سامانه، پرسش‌های پژوهشی زیر مطرح می‌گردند:
\begin{itemize}
  \item \textbf{پرسش اول ($Q_1$):} چگونه می‌توان هندسه صفحه، سرصفحه‌ها و فواصل ساختاری تعریف‌شده در شیوه‌نامه دانشگاه خوارزمی را با کمترین وابستگی به بسته‌های متداخل در بستر زی‌لاتک پیاده‌سازی نمود؟
  \item \textbf{پرسش دوم ($Q_2$):} مناسب‌ترین سازوکار برای یکپارچه‌سازی مراجع فارسی و انگلیسی بدون ایجاد خطای جهت متن و نویسه‌های گم‌شده چیست؟
  \item \textbf{پرسش سوم ($Q_3$):} طراحی ماژولار ساختار پایان‌نامه تا چه اندازه در کاهش زمان آماده‌سازی و ارتقای تجربه کاربری دانشجویان مؤثر است؟
\end{itemize}

متناظر با پرسش‌های فوق، فرضیه‌های پژوهشی زیر صورت‌بندی شده‌اند:
\begin{itemize}
  \item \textbf{فرضیه اول ($H_1$):} استفاده از ترکیب کنترل‌شده بسته‌های \lr{\texttt{geometry}}، \lr{\texttt{fancyhdr}} و \lr{\texttt{titlesec}} با تنظیمات سفارشی در کلاس پایه، امکان انطباق میلی‌متری با شیوه‌نامه دانشگاه را فراهم می‌سازد.
  \item \textbf{فرضیه دوم ($H_2$):} بهره‌گیری از سبک کتاب‌شناسی \lr{\texttt{plain-fa}} و فیلد زبانی در فایل \lr{\texttt{.bib}}، پایداری کامل ارجاعات دوزبانه را بدون هشدار خطای نویسه تضمین می‌کند.
  \item \textbf{فرضیه سوم ($H_3$):} تفکیک ساختار قالب به پوشه‌های مجزای صفحات آغازین، فصول و بخش‌های پایانی، نرخ خطای کاربران را به حداقل رسانده و یادگیری الگوها را تسهیل می‌نماید.
\end{itemize}

\section{ضرورت و اهمیت پژوهش}
\label{sec:ch1_importance}

اهمیت این پژوهش را می‌توان از دو بعد دانشگاهی و فنی تبیین نمود. از دیدگاه دانشگاهی، یکدستی و انضباط بصری رساله‌ها نشان‌دهنده هویت حرفه‌ای و انضباط سازمانی دانشگاه خوارزمی است. دیجیتالی‌سازی و انتشار الکترونیکی پایان‌نامه‌ها در کتابخانه‌های دیجیتال نیازمند فایل‌های \lr{PDF} استاندارد با برچسب‌های متنی خوانا، فراداده‌های ساخت‌یافته و پیوندهای فرامتنی پایدار است.

از بعد فنی و کاربردی، وجود این قالب به دانشجویان امکان می‌دهد به‌جای صرف صدها ساعت زمان برای برطرف‌کردن خطاهای ظاهری واژه‌پردازها، با آسودگی خاطر بر نوآوری علمی و تحلیل داده‌های پژوهش خود تمرکز نمایند. همچنین این قالب به عنوان یک الگوی بازتولیدپذیر\LTRfootnote{\lr{Reproducible Research}} برای نگارش مقالات علمی، گزارش‌های پژوهشی و کتب دانشگاهی در آینده قابل استفاده خواهد بود.

\section{ساختار پایان‌نامه}
\label{sec:ch1_organization}

ساختار این پایان‌نامه در قالب ۵ فصل و پیوست‌های مربوطه به شرح زیر سازمان‌دهی شده است:
\begin{itemize}
  \item \textbf{فصل اول (کلیات پژوهش):} به معرفی زمینه تحقیق، بیان مسئله، اهداف، پرسش‌ها، فرضیه‌ها، و ضرورت پژوهش اختصاص دارد.
  \item \textbf{فصل دوم (پیشینه پژوهش و مبانی نظری):} مبانی نظری حروف‌چینی دیجیتال، نظریه ارتباطات، مدل‌های یادگیری عمیق، و پیشینه کارهای پیشین داخلی و بین‌المللی را به تفصیل مرور و مقایسه می‌کند.
  \item \textbf{فصل سوم (روش‌شناسی و مدل پیشنهادی):} به تشریح معماری سامانه، فرمول‌بندی روابط ریاضی، ساختار الگوریتم‌ها و نحوه پیاده‌سازی اجزا می‌پردازد.
  \item \textbf{فصل چهارم (یافته‌ها و تحلیل نتایج):} نتایج تجربی، جداول ارزیابی عملکرد، نمودارهای مقایسه‌ای، شکل‌ها، قطعه کدها و تحلیل‌های کمی و کیفی را ارائه می‌دهد.
  \item \textbf{فصل پنجم (نتیجه‌گیری و پیشنهادات پژوهشی):} جمع‌بندی نهایی، دستاوردها، محدودیت‌های کار، و پیشنهادهایی برای پژوهش‌های آینده را تبیین می‌نماید.
  \item \textbf{بخش‌های پایانی:} شامل پیوست‌ها، واژه‌نامه‌های توصیفی و دوزبانه، فهرست مقالات، کارنامک علمی نویسنده، و مشخصات انگلیسی پایان‌نامه است.
\end{itemize}
       % فصل ۱: کلیات پژوهش
%% ============================================================================
%%  chapters/chapter2.tex
%%  فصل ۲: پیشینه پژوهش و مبانی نظری
%%  قالب رسمی پایان‌نامه و رساله دانشگاه خوارزمی
%% ============================================================================

\chapter{پیشینه پژوهش و مبانی نظری}
\label{chap:literature}

\section{مقدمه}
\label{sec:ch2_intro}

درک عمیق مسائل مطرح در حروف‌چینی دیجیتال، بازنمایی ساختار اسناد، و خودکارسازی صفحه‌آرایی نیازمند بررسی پایه‌های نظری در نظریه ارتباطات، سامانه‌های محاسباتی، تایپوگرافی رایانه‌ای، و رویکردهای نوین یادگیری ماشین است. در این فصل، ابتدا مبانی نظری حاکم بر پردازش اطلاعات، تکوین موتورهای حروف‌چینی دوسویه، استانداردهای بین‌المللی نویسه‌ها، و الگوهای یادگیری عمیق مورد واکاوی قرار می‌گیرد. سپس، پژوهش‌های پیشین انجام‌شده در سطح بین‌المللی و مطالعات پیشگام پژوهشگران داخلی در حوزه صفحه‌آرایی زبان فارسی و پردازش زبان طبیعی بررسی و در قالب یک جدول تطبیقی جامع مقایسه می‌گردند.

\section{مبانی نظری تحقیق}
\label{sec:ch2_theory}

\subsection{نظریه اطلاعات و بنیان‌های ارتباطی}
پایه‌های بنیادین انتقال و کدگذاری داده‌ها در سده بیستم توسط کلود شانون\LTRfootnote{\lr{Claude Elwood Shannon}} پی‌ریزی گردید. شانون در اثر تاریخی خود با عنوان نظریه ریاضی ارتباطات\index{نظریه اطلاعات} \cite{shannon1948mathematical}، مفهوم آنتروپی اطلاعات\index{آنتروپی اطلاعات} را معرفی کرد و نشان داد که محتوای پیام را می‌توان به صورت کمّی اندازه‌گیری و بدون اتلاف در کانال‌های دارای نوفه منتقل نمود. از سوی دیگر، آلن تورینگ\LTRfootnote{\lr{Alan Mathison Turing}} با طرح پرسش اساسی در خصوص قابلیت تفکر ماشین‌ها \cite{turing1950computing}، افق‌های نوینی را در پردازش نمادین و منطق محاسباتی گشود که بعدها به ظهور هوش مصنوعی\index{هوش مصنوعی} مدرن انجامید. این دو دیدگاه زیربنای نظری مدل‌سازی نمادها و متون دیجیتال را شکل می‌دهند.

\subsection{تک، لاتک و تکوین تایپوگرافی دیجیتال}
تا پیش از اواخر دهه ۱۹۷۰ میلادی، چاپ فرمول‌های ریاضی و متون علمی نیازمند فرآیندهای دستی پرهزینه و پرخطا در چاپخانه‌ها بود. دانلد کنوت\LTRfootnote{\lr{Donald Ervin Knuth}} با هدف ایجاد سامانه حروف‌چینی دقیق و مستقل از دستگاه، سامانه تک را ابداع نمود و اصول آن را در کتاب مشهور تک‌بوک \cite{knuth1984texbook} به تفصیل مدون کرد. نوآوری کنوت در ابداع الگوریتم شکست پاراگراف به سطرها بر اساس مفهوم کمینه‌سازی جریمه\LTRfootnote{\lr{Penalty and Badness Optimization}}، هنوز هم استاندارد طلایی صفحه‌آرایی متون در جهان به شمار می‌آید. پس از آن، لزلی لمپورت با توسعه بسته ماکروی لاتک، تمرکز نویسنده را از جزئیات مکانیکی صفحه به ساختار منطقی سند معطوف کرد.

\subsection{الگوریتم دوسویه یونیکد و موتور زی‌پرشین}
خط فارسی برخلاف خطوط لاتین از راست به چپ نوشته می‌شود، در حالی که ارقام، علائم جبری و عبارات لاتین تعبیه‌شده در آن جهت چپ‌به‌راست دارند. کنسرسیوم یونیکد \cite{unicode2024standard} با تدوین پیوست استاندارد شماره ۹ (\lr{UAX \#9})، الگوریتم حل جهت دوسویه\index{متن دوسویه}\LTRfootnote{\lr{Unicode Bidirectional Algorithm (Bidi)}} را معرفی نمود. وفا خلیقی \cite{khalighi2020xepersian} با تکیه بر قابلیت‌های موتور مدرن زی‌تک، بسته زی‌پرشین را پیاده‌سازی کرد. این بسته با بازتعریف پارادایم‌های پایه لاتک، پشتیبانی بی‌نقصی از ویژگی‌های ذاتی خط و زبان فارسی نظیر اتصال حروف، تنوع قلم‌های نستعلیق و نسخ، کشیدگی کلمات و تنظیم خودکار پانویس‌ها را به ارمغان آورد.

\subsection{یادگیری ماشین و شبکه‌های عصبی عمیق}
در دهه گذشته، پارادایم پردازش اسناد و تحلیل متون با پیشرفت‌های شگرف یادگیری ماشین\index{یادگیری ماشین} دگرگون شده است. مبانی یادگیری آماری و شناسایی الگو توسط بیشف \cite{bishop2006pattern} و اصول یادگیری عمیق\index{یادگیری عمیق} توسط گودفلو و همکاران \cite{goodfellow2016deep} تدوین شده است. در جامعه علمی کشور نیز کتاب مبانی یادگیری ماشین کریمی \cite{karimi1398} و ترجمه فارسی مرجع جامع هوش مصنوعی راسل و نورویگ توسط سلیمی \cite{salimi1400translation} به تبیین دقیق این مفاهیم برای دانشجویان و پژوهشگران پرداخته‌اند.

علاوه بر این، معرفی معماری ترنسفورمر\index{ترنسفورمر} و سازوکار توجه خودکار\index{سازوکار توجه} توسط واسوانی و همکاران \cite{vaswani2017attention}، انقلابی در پردازش توالی‌ها پدید آورد که منجر به پیدایش مدل‌های زبانی عظیم گردید. در حوزه پردازش تصویر و بازشناسی نویسه‌ها، شبکه‌های زاینده تخاصمی\index{شبکه زاینده تخاصمی} گودفلو و همکاران \cite{goodfellow2014generative} و شبکه‌های عمیق بازمانده\index{شبکه عصبی پیچشی} هه و همکاران \cite{he2016deep} بن‌بست‌های یادگیری در شبکه‌های بسیار عمیق را مرتفع ساختند.

\section{مرور پژوهش‌های پیشین}
\label{sec:ch2_related_works}

\subsection{پژوهش‌های بین‌المللی در خودکارسازی صفحه‌آرایی اسناد}
تحقیقات در زمینه درک ساختار اسناد\LTRfootnote{\lr{Document Layout Analysis}} و تولید خودکار متون دانشگاهی طی سال‌های اخیر شتاب فزاینده‌ای یافته است. محققان دانشگاه‌های استنفورد و ام‌آی‌تی با توسعه بسته‌های متعددی برای مجلات معتبر \lr{IEEE} و \lr{ACM}، استانداردسازی خروجی مقالات را محقق ساختند. تمرکز اصلی این پژوهش‌ها بر تفکیک محتوا از قالب، کاهش بار محاسباتی کامپایلر، و تولید اسناد دسترسی‌پذیر\LTRfootnote{\lr{Accessible PDF}} بوده است. با این حال، غالب این سامانه‌ها برای زبان‌های تک‌جهته طراحی شده و در تعامل با خطوط عربی و فارسی با محدودیت‌های اساسی روبه‌رو هستند.

\subsection{پژوهش‌های داخلی در نگارش متون دانشگاهی و پردازش خط فارسی}
در داخل کشور، تلاش‌های متعددی برای استانداردسازی حروف‌چینی و پردازش متون علمی صورت پذیرفته است:
\begin{itemize}
  \item \textbf{مطالعه سامانه‌های اطلاعاتی و رفتار پژوهشگران:} رجبعلی‌بگلو، فتاحی و پریرخ \cite{fatahi1395} به بررسی تأثیر ساختار نظام‌های اطلاعاتی و قالب‌های نگارشی بر رفتار اطلاع‌یابی و بهره‌وری پژوهشگران پرداختند و نشان دادند که استانداردسازی فرمت اسناد، سرعت بازیابی اطلاعات و ارجاع‌دهی را به میزان چشمگیری افزایش می‌دهد.
  \item \textbf{مدل‌سازی زبان فارسی با شبکه‌های عمیق:} دهقان \cite{dehghan1399phd} در رساله دکترای خود در دانشگاه خوارزمی به بررسی بهینه‌سازی شبکه‌های عصبی بازگشتی و ترنسفورمرها در پردازش زبان طبیعی فارسی پرداخت و چالش‌های واژه‌نگاشت، فاصله‌گذاری مجازی و ساختار دستوری متون علمی فارسی را مورد ارزیابی قرار داد.
  \item \textbf{بهینه‌سازی توزیع‌شده در یادگیری عمیق:} شمس، محمدی و رضایی \cite{shams1401} الگوریتم‌های بهینه‌سازی توزیع‌شده را برای آموزش مدل‌های عمیق بررسی نمودند که نتایج آن در تسریع زمان آموزش سامانه‌های هوشمند پردازش سند بسیار راه‌گشا است.
  \item \textbf{صفحه‌آرایی خودکار در زبان‌های راست‌به‌چپ:} احمدی \cite{ahmadi1401thesis} در پایان‌نامه کارشناسی ارشد خود در دانشکده مهندسی برق و کامپیوتر دانشگاه خوارزمی، مسئله طراحی سامانه‌های خودکار صفحه‌آرایی اسناد دانشگاهی را در بستر زبان‌های راست‌به‌چپ تحلیل و ساختار کلاس‌های تک را بررسی نمود.
  \item \textbf{تحلیل مقایسه‌ای سامانه‌های حروف‌چینی متن‌باز:} رضایی و حسینی \cite{rezaei1402} در کنفرانس ملی کامپیوتر ایران، عملکرد سامانه‌های مختلف حروف‌چینی را در تولید مستندات علمی دوزبانه مقایسه کردند و نشان دادند که زی‌پرشین در مقایسه با سایر ابزارها دارای بالاترین امتیاز در حفظ زیبایی بصری و انطباق با ضوابط دانشگاهی است.
\end{itemize}

\section{تحلیل مقایسه‌ای و جدول تطبیقی کارهای مرتبط}
\label{sec:ch2_comparison}

به منظور ارزیابی جامع رویکردهای موجود و تبیین وجوه تمایز قالب ارائه‌شده در این پژوهش با سایر ابزارها، جدول \ref{tab:comparison_literature} ویژگی‌های کلیدی سامانه‌های متداول را مقایسه می‌نماید.

\begin{table}[htbp]
\centering
\caption{مقایسه تطبیقی سامانه‌های حروف‌چینی و قالب‌های اسناد دانشگاهی}
\label{tab:comparison_literature}
\vspace{4pt}
\small
\setlength{\tabcolsep}{4pt}
\resizebox{\linewidth}{!}{%
\begin{tabular}{|p{2.3cm}|c|c|c|c|p{4.2cm}|}
\hline
\textbf{سامانه / رویکرد} & \textbf{محیط اجرا} & \textbf{پشتیبانی دوزبانه} & \textbf{پایداری ارجاعات} & \textbf{انطباق با خوارزمی} & \textbf{مهم‌ترین مزایا و محدودیت‌ها} \\
\hline
مایکروسافت ورد & تجاری & متوسط & ضعیف & اسمی (فایل الگو) & سهولت کاربری اولیه، اما به‌هم‌ریختگی در اسناد حجیم و فرمول‌ها \\
\hline
لاتک استاندارد & متن‌باز & نامناسب (نیازمند عرب‌تک) & عالی & ندارد & توانایی بالا در متون انگلیسی، ناتوانی در قلم‌های بومی فارسی \\
\hline
قالب‌های متفرقه وب & متن‌باز & ناپایدار & متوسط & ندارد & عدم هماهنگی با حاشیه‌ها و سرصفحه‌های شیوه‌نامه مصوب ۱۴۰۲ \\
\hline
قالب دانشگاه تهران & متن‌باز & خوب & خوب & مغایر در حاشیه‌ها & تنظیم حاشیه‌ها بر اساس ۳/۵ و ۲/۵ سانتی‌متر (ناهمخوان با خوارزمی) \\
\hline
\textbf{قالب پیشنهادی خوارزمی} & \textbf{متن‌باز} & \textbf{کامل و پایدار} & \textbf{عالی (دوزبانه)} & \textbf{۱۰۰ درصد منطبق} & \textbf{طراحی ماژولار، قلم‌های رسمی بی زر و بی تیتر، حذف کامل هشدارهای نویسه} \\
\hline
\end{tabular}%
}
\end{table}

\section{جمع‌بندی فصل}
\label{sec:ch2_summary}

بررسی مبانی نظری و پیشینه پژوهش‌های پیشین نشان می‌دهد که حروف‌چینی اسناد علمی فارسی در سطوح عالی دانشگاهی، نیازمند اتخاذ رویکردی چندبعدی شامل تسلط بر استانداردهای یونیکد، قواعد دوسویه، معماری ماژولار لاتک و تطبیق ریاضی با فواصل شیوه‌نامه دانشگاه خوارزمی است. اگرچه تلاش‌های ارزنده‌ای توسط پژوهشگران گذشته صورت گرفته، اما خلأ یک قالب یکپارچه، کاملاً آزموده‌شده، فاقد هرگونه خطای نویسه گم‌شده، و همراه با مستندات کامل برای دانشجویان دانشگاه خوارزمی همچنان احساس می‌شد. مدل پیشنهادی که در فصل بعدی تشریح می‌گردد، پاسخگوی دقیق این نیازمندی‌ها خواهد بود.
       % فصل ۲: پیشینه پژوهش و مبانی نظری
%% ============================================================================
%%  chapters/chapter3.tex
%%  فصل ۳: روش‌شناسی و مدل پیشنهادی
%%  قالب رسمی پایان‌نامه و رساله دانشگاه خوارزمی
%% ============================================================================

\chapter{روش‌شناسی و مدل پیشنهادی}
\label{chap:methodology}

\section{مقدمه}
\label{sec:ch3_intro}

دستیابی به یک ساختار حروف‌چینی پایدار و خودکار برای پایان‌نامه‌ها و رساله‌های دانشگاهی، نیازمند طراحی یک متدولوژی علمی مدون و معماری چندلایه‌ای است که بتواند پیچیدگی‌های مربوط به فونت‌های فارسی، قوانین چیدمان دوسویه، شماره‌گذاری سلسله‌مراتبی، و فرمول‌بندی‌های پیشرفته ریاضی را مدیریت کند. در این فصل، ابتدا معماری کلان سامانه پیشنهادی و تفکیک لایه‌های کارکردی آن تشریح می‌گردد. سپس مدل‌سازی و فرمول‌بندی ریاضی اجزای سامانه شامل توابع هدف، روابط بهینه‌سازی ماتریسی، توابع تکه‌ای و انتگرال‌های احتمالاتی تبیین می‌شود. در ادامه، الگوریتم جامع پردازش و حل بهینه‌سازی بر پایه ساختار الگوریتمی ارائه‌شده و مشخصات فنی پیاده‌سازی ماژولار قالب به تفصیل بیان خواهد شد.

\section{معماری کلی سامانه و ساختار ماژولار}
\label{sec:ch3_architecture}

معماری سامانه پیشنهادی بر پایه اصل تفکیک دغدغه‌ها\LTRfootnote{\lr{Separation of Concerns (SoC)}} بنا نهاده شده است. بر این اساس، ساختار کلی پایان‌نامه به ۵ لایه مجزا تقسیم می‌گردد:
\begin{enumerate}
  \item \textbf{لایه هسته ساختاری و سندپایه (\lr{\texttt{kharazmi-thesis.cls}}):} این لایه وظیفه بارگذاری کلاس پایه گزارش، اعمال دقیق ابعاد هندسی کاغذ \lr{A4}، تعریف حاشیه‌ها (بالا ۳/۵، پایین ۲/۵، راست ۳/۰ و چپ ۳/۰ سانتی‌متر)، فاصله سرصفحه‌ها و پاصفحه‌ها، قالب‌بندی سلسله‌مراتب عناوین فصل‌ها و بخش‌ها، و قالب‌بندی زیرنویس جدول‌ها و شکل‌ها را بر عهده دارد.
  \item \textbf{لایه استایل و فراداده (\lr{\texttt{kharazmi.sty}}):} در این لایه، تنظیمات قلم‌های رسمی فارسی (بی زر ۱۳ برای متن و بی تیتر برای سرعنوان‌ها)، مقیاس‌بندی دقیق قلم لاتین تایمز نیو رومن، ماکروهای نگهداری مشخصات پایان‌نامه، و توابع تولیدکننده برگه‌های رسمی پیش از متن و پس از متن تعریف شده‌اند.
  \item \textbf{لایه صفحات آغازین (\lr{\texttt{front/}}):} شامل ماژول‌های مجزای جلد فارسی، صفحه بسم‌الله، صورت‌جلسه دفاع و تصویب‌نامه هیئت داوران، تعهدنامه اصالت اثر، تقدیم، سپاسگزاری، و چکیده فارسی است.
  \item \textbf{لایه فصول محتوایی (\lr{\texttt{chapters/}}):} ۵ فصل مستقل پایان‌نامه که به صورت فایل‌های مجزا توسط دستور \lr{\texttt{\textbackslash include}} به فایل اصلی متصل می‌شوند.
  \item \textbf{لایه مراجع و صفحات پایانی (\lr{\texttt{back/}} و \lr{\texttt{references.bib}}):} دربرگیرنده مراجع دوزبانه، پیوست‌های ریاضی و اثبات‌ها، واژه‌نامه‌های سه‌گانه، کارنامک علمی، و برگه‌های انگلیسی چکیده و جلد پشتی است.
\end{enumerate}

\section{فرمول‌بندی ریاضی و مدل‌سازی تحلیلی}
\label{sec:ch3_formulation}

در این بخش، به منظور ارائه نمونه‌ای کامل از قابلیت‌های پیشرفته ریاضی قالب دانشگاه خوارزمی، انواع گوناگون روابط ریاضی شامل معادلات تک‌خطی، چندخطی هم‌تراز، ماتریس‌ها، توابع تکه‌ای، و انتگرال‌ها فرمول‌بندی می‌گردد. تمامی شماره‌گذاری‌های روابط با استفاده از شناسه فصل و شماره رابطه به صورت فارسی (مانند \eqref{eq:ch3_single}) حروف‌چینی می‌شوند.

\subsection{معادلات تک‌خطی و تابع هدف}
فرض کنید $\mathcal{D} = \{(x_i, y_i)\}_{i=1}^{N}$ مجموعه داده‌های آموزشی شامل $N$ نمونه مستقل باشد. تابع هدف بهینه‌سازی\index{بهینه‌سازی} پارامترهای مدل $\theta \in \mathbb{R}^d$ با در نظر گرفتن ترم تنظیم‌سازی\index{تنظیم‌سازی} تیخونوف\LTRfootnote{\lr{Tikhonov Regularization}} به صورت رابطه تک‌خطی \eqref{eq:ch3_single} تعریف می‌شود:
\begin{equation}
\label{eq:ch3_single}
\mathcal{J}(\theta) = \frac{1}{N} \sum_{i=1}^{N} \ell\left(f(x_i; \theta), y_i\right) + \frac{\lambda}{2} \|\theta\|_2^2
\end{equation}
که در آن $\ell(\cdot, \cdot)$ نشان‌دهنده تابع زیان\index{تابع زیان} محلی بر روی هر نمونه، $\|\theta\|_2 = \sqrt{\sum_{j=1}^d \theta_j^2}$ بیانگر نُرم اقلیدسی بردار پارامترها، و $\lambda > 0$ ضریب کنترل پیچیدگی مدل است.

\subsection{معادلات چندخطی هم‌تراز}
برای به‌روزرسانی پارامترها در فرآیند گرادیان کاهشی\index{گرادیان کاهشی} تصادفی با تکانه مومنتوم نسترُف\LTRfootnote{\lr{Nesterov Accelerated Gradient}}، دنباله محاسبات به صورت رابطه چندخطی هم‌تراز \eqref{eq:ch3_multiline} صورت می‌پذیرد:
\begin{align}
\label{eq:ch3_multiline}
v_{t+1} &= \mu \, v_t - \eta \, \nabla_{\theta} \mathcal{J}\left(\theta_t + \mu \, v_t\right) \\
\theta_{t+1} &= \theta_t + v_{t+1} \\
\Delta \theta_t &= \|\theta_{t+1} - \theta_t\|_2 = \|v_{t+1}\|_2
\end{align}
که در آن $\mu \in [0, 1)$ ضریب تکانه، $\eta > 0$ نرخ یادگیری، و $v_t$ بردار سرعت در تکرار $t$ است.

\subsection{مدل‌سازی ماتریسی و سازوکار توجه خودکار}
در مدل‌های مدرن مبتنی بر توجه خودکار، ماتریس‌های پرسش ($\mathbf{Q}$)، کلید ($\mathbf{K}$) و مقدار ($\mathbf{V}$) از ضرب ماتریس ویژگی‌های ورودی $\mathbf{X} \in \mathbb{R}^{n \times d}$ در ماتریس‌های وزن پدید می‌آیند:
\begin{equation}
\label{eq:ch3_matrices}
\begin{bmatrix}
\mathbf{Q} \\
\mathbf{K} \\
\mathbf{V}
\end{bmatrix}
=
\begin{bmatrix}
\mathbf{W}_Q \\
\mathbf{W}_K \\
\mathbf{W}_V
\end{bmatrix}
\mathbf{X}^\top
\end{equation}
خروجی لایه توجه مقیاس‌شده به صورت ضرب ماتریسی در رابطه \eqref{eq:ch3_attention} مدل‌سازی می‌گردد:
\begin{equation}
\label{eq:ch3_attention}
\operatorname{Attention}(\mathbf{Q}, \mathbf{K}, \mathbf{V}) = \operatorname{softmax}\left( \frac{\mathbf{Q} \mathbf{K}^\top}{\sqrt{d_k}} \right) \mathbf{V}
\end{equation}
به طوری که ماتریس هسته توجه $\mathbf{A} \in \mathbb{R}^{n \times n}$ دارای درایه‌هایی به صورت زیر است:
\begin{equation}
\label{eq:ch3_softmax_matrix}
A_{ij} = \frac{\exp\left(\frac{q_i^\top k_j}{\sqrt{d_k}}\right)}{\sum_{l=1}^n \exp\left(\frac{q_i^\top k_l}{\sqrt{d_k}}\right)}
\end{equation}

\subsection{توابع چندضابطه‌ای (تکه‌ای)}
توابع فعال‌ساز غیرخطی در لایه‌های پنهان شبکه‌های عصبی نقش اساسی دارند. تابع فعال‌ساز اصلاح‌شده تکه‌ای خطی\LTRfootnote{\lr{Parametric Rectified Linear Unit (PReLU)}} به صورت رابطه \eqref{eq:ch3_piecewise} تعریف می‌شود:
\begin{equation}
\label{eq:ch3_piecewise}
\phi(z) =
\begin{cases}
z & \text{اگر } z \ge 0 \\
\alpha z & \text{اگر } z < 0
\end{cases}
\end{equation}
همچنین تابع زیان هابر\LTRfootnote{\lr{Huber Loss Function}} که نسبت به داده‌های پرت مقاوم است، ضابطه‌ای دوشاخه دارد:
\begin{equation}
\label{eq:ch3_huber}
\mathcal{L}_\delta(e) =
\begin{cases}
\frac{1}{2} e^2 & \text{اگر } |e| \le \delta \\[1em]
\delta \left( |e| - \frac{1}{2} \delta \right) & \text{وگرنه}
\end{cases}
\end{equation}

\subsection{روابط انتگرالی و محاسبات پیوسته}
در تحلیل‌های احتمالاتی و برآورد درست‌نمایی بیزی، انتگرال‌های چندبعدی بر روی توزیع پیشین پارامترها به کار می‌روند. امید ریاضی تابع زیان روی فضای پیوسته به صورت رابطه \eqref{eq:ch3_integral} محاسبه می‌گردد:
\begin{equation}
\label{eq:ch3_integral}
\mathbb{E}_{x \sim p(x)}\left[\mathcal{L}(x)\right] = \int_{\mathcal{X}} \mathcal{L}(x) \, p(x) \, \mathrm{d}x = \int_{-\infty}^{+\infty} \cdots \int_{-\infty}^{+\infty} \mathcal{L}(x_1, \dots, x_d) \, p(x_1, \dots, x_d) \, \mathrm{d}x_1 \cdots \mathrm{d}x_d
\end{equation}
همچنین احتمال پسین به وسیله قاعده بیز به صورت زیر فرمول‌بندی می‌شود:
\begin{equation}
\label{eq:ch3_bayes}
p(\theta \mid \mathcal{D}) = \frac{p(\mathcal{D} \mid \theta) \, p(\theta)}{\int_{\Theta} p(\mathcal{D} \mid \theta') \, p(\theta') \, \mathrm{d}\theta'}
\end{equation}

\section{طراحی الگوریتم پیشنهادی}
\label{sec:ch3_algorithm}

الگوریتم جامع بهینه‌سازی ساختار سند و آموزش مدل پیشنهادی در الگوریتم \ref{alg:ch3_optimization} با کلمات کلیدی فارسی به تصویر کشیده شده است. این الگوریتم با دریافت داده‌ها و فراسنج‌های اولیه، در هر دور فرآیند گرادیان کاهشی را اعمال کرده و شرایط توقف زودهنگام را می‌آزماید.

\begin{algorithm}[htbp]
\caption{آموزش و بهینه‌سازی پارامترهای مدل پیشنهادی}
\label{alg:ch3_optimization}
\begin{algorithmic}[1]
\REQUIRE مجموعه داده آموزشی $\mathcal{D}$، نرخ یادگیری $\eta$، تعداد کل دورها $E$، آستانه همگرایی $\epsilon$
\ENSURE بردار بهینه پارامترها $\theta^*$ و تاریخچه تابع زیان $\mathcal{H}$
\STATE بردار پارامترها را به صورت تصادفی مقداردهی اولیه کن: $\theta \leftarrow \theta_0$
\STATE متغیر تاریخچه را خالی در نظر بگیر: $\mathcal{H} \leftarrow \emptyset$
\FOR{$e = 1$ \TO $E$}
  \STATE داده‌های آموزشی $\mathcal{D}$ را به صورت تصادفی بر بزن
  \FOR{هر زیردسته داده $\mathcal{B} \subset \mathcal{D}$ با اندازه $B$}
    \STATE گرادیان تابع زیان روی زیردسته را محاسبه کن: $g \leftarrow \frac{1}{B} \sum_{x \in \mathcal{B}} \nabla_\theta \ell(x; \theta)$
    \STATE پارامترها را به‌روزرسانی کن: $\theta \leftarrow \theta - \eta \cdot g$
  \ENDFOR
  \STATE مقدار کل زیان را در پایان دور محاسبه کن: $L_e \leftarrow \mathcal{J}(\theta)$
  \STATE مقدار جدید را به تاریخچه بیفزا: $\mathcal{H} \leftarrow \mathcal{H} \cup \{L_e\}$
  \IF{$e > 1$ \AND $|L_e - L_{e-1}| < \epsilon$}
    \PRINT پیام همگرایی پیش از موعد در دور $e$
    \RETURN $\theta$
  \ENDIF
\ENDFOR
\RETURN $\theta$
\end{algorithmic}
\end{algorithm}

\section{پیاده‌سازی ماژولار و رابط‌های نرم‌افزاری}
\label{sec:ch3_interfaces}

به منظور اطمینان از سازگاری و انعطاف‌پذیری قالب، اجزای مختلف از طریق واسط‌های استاندارد با یکدیگر تعامل دارند. هر بخش از مستند دارای پرونده مستقل است و کاربر بدون نیاز به باز کردن یا دستکاری کدهای بنیادین کلاس، تنها با ویرایش فراداده‌های موجود در پرونده \lr{\texttt{main.tex}} اطلاعات شخصی، اسامی اساتید، عنوان و تاریخ را مقداردهی می‌کند. ساختار کلاس به گونه‌ای محافظت شده است که در صورت افزودن بسته‌های تکمیلی، هماهنگی سرصفحه‌ها و حاشیه‌های مصوب دانشگاه خوارزمی مخدوش نشود.

\section{جمع‌بندی فصل}
\label{sec:ch3_summary}

در این فصل، معماری ۵ لایه سامانه حروف‌چینی اسناد دانشگاهی تشریح و فرمول‌بندی تحلیلی انواع روابط ریاضی شامل معادلات تک‌خطی، چندخطی هم‌تراز، ماتریس‌ها، توابع تکه‌ای، و انتگرال‌های پیوسته به همراه شماره‌گذاری استاندارد فارسی ارائه گردید. همچنین ساختار الگوریتمی آموزش و بهینه‌سازی با کلمات کلیدی بومی‌سازی‌شده تدوین شد. در فصل چهارم، نتایج حاصل از ارزیابی عملی این مدل، جداول آماری، نمودارها، شکل‌ها و قطعه کدها به تفصیل تحلیل خواهند شد.
       % فصل ۳: روش‌شناسی و مدل پیشنهادی
%% ============================================================================
%%  chapters/chapter4.tex
%%  فصل ۴: یافته‌ها و تحلیل نتایج
%%  قالب رسمی پایان‌نامه و رساله دانشگاه خوارزمی
%% ============================================================================

\chapter{یافته‌ها و تحلیل نتایج}
\label{chap:results}

\section{مقدمه}
\label{sec:ch4_intro}

ارزیابی تجربی و اعتبارسنجی عملی هر سامانه نرم‌افزاری و الگوریتم محاسباتی، بخش اساسی در اثبات کارایی، پایداری و صحت آن به شمار می‌رود. در این فصل، نتایج پیاده‌سازی قالب جامع پایان‌نامه دانشگاه خوارزمی و مدل‌های پیشنهادی مورد ارزیابی دقیق کمّی و کیفی قرار می‌گیرد. ابتدا مشخصات محیط سخت‌افزاری و نرم‌افزاری آزمایشگاه و ویژگی‌های دادگان مورد بررسی تشریح می‌گردد. سپس معیارهای ارزیابی تعریف شده و جداول مقایسه‌ای عملکرد با استفاده از امکانات پیشرفته جدول‌بندی لاتک ارائه می‌شوند. در ادامه، نمودارهای روند همگرایی، تصویرهای خروجی به همراه زیرشکل‌ها، و قطعه‌کدهای پیاده‌سازی گزارش شده و در نهایت، تحلیل حساسیت و بحث پیرامون یافته‌ها صورت می‌پذیرد.

\section{محیط آزمایشگاهی و مجموعه‌داده‌های پژوهش}
\label{sec:ch4_setup}

کلیه آزمایش‌ها و فرآیندهای کامپایل چندمرحله‌ای بر روی یک سامانه محاسباتی با پردازنده \lr{Intel Core i9-13900K} با فرکانس کاری ۵/۴ گیگاهرتز، ۶۴ گیگابایت حافظه رم از نوع \lr{DDR5}، و شتاب‌دهنده گرافیکی \lr{NVIDIA GeForce RTX 4090} با ۲۴ گیگابایت حافظه اختصاصی اجرا شده است. بستر نرم‌افزاری شامل توزیع میک‌تک نسخه ۲۶/۵ بر روی سیستم‌عامل ویندوز، موتور حروف‌چینی زی‌تک ویرایش ۰/۹۹۹۹۹۸، پایتون ۳/۱۰ و کتابخانه‌های یادگیری ماشین پای‌تورچ ۲/۲ است.

مجموعه‌داده مورد استفاده در ارزیابی سامانه شامل ۵۰ عنوان رساله و پایان‌نامه واقعی دانشکده مهندسی برق و کامپیوتر دانشگاه خوارزمی در حوزه‌های هوش مصنوعی، مخابرات و نرم‌افزار است که ساختارهای پیچیده‌ای از متون دوزبانه، فرمول‌های ریاضی چندسطحی، جداول بزرگ و اشکال با وضوح بالا را در بر دارند.

\section{معیارهای ارزیابی کارایی}
\label{sec:ch4_metrics}
\index{معیارهای ارزیابی}

به منظور ارزیابی دقیق کیفیت ساختاری و سرعت محاسباتی سامانه، معیارهای آماری زیر به کار گرفته شده‌اند:
\begin{enumerate}
  \item \textbf{دقت (\lr{Precision}):} نسبت عناصر ساختاری شناسایی‌شده درست به کل عناصر شناسایی‌شده:
  \begin{equation}
  \label{eq:ch4_precision}
  P = \frac{TP}{TP + FP}
  \end{equation}
  \item \textbf{بازخوانی (\lr{Recall}):} نسبت عناصر ساختاری درست شناسایی‌شده به کل عناصر واقعی موجود در سند:
  \begin{equation}
  \label{eq:ch4_recall}
  R = \frac{TP}{TP + FN}
  \end{equation}
  \item \textbf{امتیاز اف‌یک (\lr{F1-Score}):} میانگین همساز دقت و بازخوانی جهت ارزیابی متوازن عملکرد:
  \begin{equation}
  \label{eq:ch4_f1}
  F_1 = 2 \times \frac{P \times R}{P + R}
  \end{equation}
  \item \textbf{زمان کامپایل (\lr{Compilation Time}):} کل زمان لازم برای کامپایل کامل چندمرحله‌ای سند (شامل اجرای زی‌لاتک، بیب‌تک، و دو بار اجرای مجدد زی‌لاتک جهت تثبیت شماره‌گذاری‌ها و پیوندها) بر حسب ثانیه.
  \item \textbf{بیشینه مصرف حافظه (\lr{Peak Memory Footprint}):} حداکثر حجم حافظه اصلی اشغال‌شده در طول فرآیند پردازش و حروف‌چینی بر حسب مگابایت.
\end{enumerate}

\section{تحلیل نتایج تجربی و جداول مقایسه‌ای}
\label{sec:ch4_tables}
\index{ماتریس درهم‌ریختگی}

در جدول \ref{tab:model_performance}، مقایسه عملکرد روش پایه خطی، شبکه عصبی چندلایه (\lr{MLP})، شبکه عصبی پیچشی (\lr{CNN}) و مدل پیشنهادی ترنسفورمر ماژولار آورده شده است. این جدول با بهره‌گیری از بسته‌های حرفه‌ای \lr{\texttt{booktabs}} و \lr{\texttt{multirow}} حروف‌چینی شده است.

\begin{table}[htbp]
\centering
\caption{ارزیابی مقایسه‌ای عملکرد مدل‌های مختلف در بازشناسی ساختار سند}
\label{tab:model_performance}
\vspace{4pt}
\begin{tabular}{lcccccc}
\toprule
\textbf{مدل ارزیابی‌شده} & \textbf{دقت (\%)} & \textbf{بازخوانی (\%)} & \textbf{امتیاز $F_1$ (\%)} & \textbf{زمان کامپایل (\lr{s})} & \textbf{حافظه (\lr{MB})} \\
\midrule
مدل پایه رگرسیون خطی & ۸۲/۴ & ۷۹/۱ & ۸۰/۷ & ۱/۴ & ۱۲۸ \\
شبکه چندلایه پرسپترون & ۸۸/۷ & ۸۶/۳ & ۸۷/۵ & ۲/۸ & ۲۵۶ \\
شبکه عصبی پیچشی عمیق & ۹۳/۲ & ۹۱/۸ & ۹۲/۵ & ۴/۵ & ۵۱۲ \\
\textbf{مدل پیشنهادی ماژولار} & \textbf{۹۸/۶} & \textbf{۹۷/۹} & \textbf{۹۸/۲} & \textbf{۳/۲} & \textbf{۳۸۴} \\
\bottomrule
\end{tabular}
\end{table}

همچنین جدول \ref{tab:detailed_breakdown} جزئیات نرخ خطای قالب را در بخش‌های مختلف ساختاری شامل حاشیه‌ها، سرصفحه‌ها، پانویس‌ها، فرمول‌ها و مراجع دوزبانه نشان می‌دهد.

\begin{table}[htbp]
\centering
\caption{ارزیابی انطباق اجزای قالب با شیوه‌نامه دانشگاه خوارزمی}
\label{tab:detailed_breakdown}
\vspace{4pt}
\begin{tabular}{llccp{5cm}}
\toprule
\textbf{مؤلفه ساختاری} & \textbf{شاخص مورد سنجش} & \textbf{تعداد آزمون} & \textbf{نرخ موفقیت (\%)} & \textbf{وضعیت انطباق با شیوه‌نامه} \\
\midrule
\multirow{2}{*}{هندسه صفحه} & حاشیه‌های چهارگانه & ۵۰ & ۱۰۰/۰ & منطبق دقیق بر ۳/۵، ۲/۵، ۳/۰، ۳/۰ \lr{cm} \\
& فاصله سرصفحه (۲/۲۵ \lr{cm}) & ۵۰ & ۱۰۰/۰ & منطبق بر خط جداکننده ۰/۶ \lr{pt} \\
\midrule
\multirow{2}{*}{تایپوگرافی} & قلم متن اصلی (بی زر ۱۳) & ۵۰ & ۱۰۰/۰ & بدون هشدار فونت با مقیاس ۱/۰۸ \\
& قلم عناوین (بی تیتر) & ۵۰ & ۱۰۰/۰ & رعایت سلسله‌مراتب ۱۶، ۱۳ و ۱۲ \lr{pt} \\
\midrule
\multirow{2}{*}{عناصر ریاضی} & فرمول‌های تک و چندخطی & ۱۲۰ & ۱۰۰/۰ & شماره‌گذاری فارسی (فصل-رابطه) \\
& نمادها و ماتریس‌ها & ۸۰ & ۱۰۰/۰ & هم‌ترازی صحیح در جهت فرمول \\
\midrule
\multirow{2}{*}{بخش‌های دوزبانه} & پانویس‌های درون‌متنی & ۳۵۰ & ۱۰۰/۰ & ریست شماره در هر صفحه \\
& مراجع فارسی و انگلیسی & ۲۰۰ & ۱۰۰/۰ & کامپایل بیب‌تک با سبک \lr{plain-fa} \\
\bottomrule
\end{tabular}
\end{table}

\section{نمودارهای ارزیابی و رفتار همگرایی}
\label{sec:ch4_charts}

در قالب پایان‌نامه دانشگاه خوارزمی، نمودارها دارای محیط مستقل شناور (\lr{\texttt{chart}}) با نام‌گذاری «نمودار» و فهرست اختصاصی هستند. نمودار \ref{chart:loss_convergence} تغییرات تابع زیان و دقت یادگیری را در طول ۱۰۰ دور متوالی آموزش نمایش می‌دهد.

\begin{chart}[htbp]
\centering
\fbox{\parbox{12cm}{%
  \vspace{1.5cm}
  \centering
  \textbf{[نمودار ۴-۱: روند کاهش تابع زیان تجربی و افزایش دقت در طول دورهای آموزش]}\par
  \vspace{0.4cm}
  \small
  دور ۱: زیان = ۲/۴۵، دقت = ۵۴/۲ درصد $\quad\longrightarrow\quad$ دور ۵۰: زیان = ۰/۲۸، دقت = ۹۴/۱ درصد $\quad\longrightarrow\quad$ دور ۱۰۰: زیان = ۰/۰۴، دقت = ۹۸/۶ درصد\par
  \vspace{1.5cm}
}}
\caption{روند همگرایی تابع زیان و ارتقای دقت مدل در طول دورهای متوالی آموزش}
\label{chart:loss_convergence}
\end{chart}

\section{شکل‌ها و زیرشکل‌های تصویری}
\label{sec:ch4_figures}

در این بخش، نمونه‌ای از درج شکل‌های تصویری تکی و مرکب ارائه می‌شود. در شکل \ref{fig:khu_compound}، نشان رسمی دانشگاه خوارزمی و ساختار بصری خروجی قالب در قالب دو زیرشکل الف و ب نمایش داده شده است. زیرنویس شکل‌ها بر اساس ضوابط دانشگاه خوارزمی با قلم بی زر ضخیم ۹ در زیر تصویر درج می‌گردد.

\begin{figure}[htbp]
\centering
\begin{subfigure}[b]{0.42\textwidth}
  \centering
  \includegraphics[height=3.2cm]{khu_logo.png}
  \caption{نشان رسمی دانشگاه خوارزمی}
  \label{subfig:khu_logo}
\end{subfigure}
\hfill
\begin{subfigure}[b]{0.52\textwidth}
  \centering
  \includegraphics[height=3.2cm]{sample_fig.png}
  \caption{نمونه ساختار بصری صفحه‌آرایی}
  \label{subfig:sample_page}
\end{subfigure}
\caption{تصاویر اعتبارسنجی بصری: (الف) نشان دانشگاه خوارزمی؛ (ب) ساختار لایه‌بندی قالب}
\label{fig:khu_compound}
\end{figure}

همچنین شکل \ref{fig:sample_architecture} نمای بزرگ‌تری از نمودار چیدمان ماژول‌های قالب را ارائه می‌کند:

\begin{figure}[htbp]
\centering
\includegraphics[width=0.75\textwidth]{sample_fig.png}
\caption{نمای شماتیک معماری ماژولار و گردش کار کامپایل اسناد در سامانه خوارزمی}
\label{fig:sample_architecture}
\end{figure}

\section{قطعه کدهای پیاده‌سازی و ارزیابی خودکار}
\label{sec:ch4_listings}

برای مستندسازی کدهای محاسباتی در متن پایان‌نامه، محیط \lr{\texttt{listings}} به همراه شماره‌گذاری سطرها و کادر مشخص تعبیه شده است. کد \ref{lst:metric_eval} تابع پایتون جهت محاسبه خودکار ماتریس درهم‌ریختگی و معیارهای دقت، بازخوانی و امتیاز اف‌یک را نشان می‌دهد.

\begin{latin}
\begin{lstlisting}[language=Python, caption={\rl{تابع پایتون جهت محاسبه ماتریس درهم‌ریختگی و معیارهای ارزیابی}}]
import numpy as np

def compute_evaluation_metrics(y_true, y_pred):
    """
    Computes precision, recall, and F1-score from confusion matrix.
    """
    classes = np.unique(y_true)
    n_classes = len(classes)
    confusion = np.zeros((n_classes, n_classes), dtype=int)
    
    for t, p in zip(y_true, y_pred):
        confusion[t, p] += 1
        
    tp = np.diag(confusion)
    fp = np.sum(confusion, axis=0) - tp
    fn = np.sum(confusion, axis=1) - tp
    
    precision = np.mean(tp / (tp + fp + 1e-12))
    recall = np.mean(tp / (tp + fn + 1e-12))
    f1 = 2 * (precision * recall) / (precision + recall + 1e-12)
    
    return {
        "precision": float(precision),
        "recall": float(recall),
        "f1_score": float(f1),
        "confusion_matrix": confusion.tolist()
    }
\end{lstlisting}
\label{lst:metric_eval}
\end{latin}

\section{بحث و تحلیل حساسیت}
\label{sec:ch4_discussion}

تحلیل‌های آماری انجام‌شده نشان می‌دهند که خطای کامپایل و هشدارهای مربوط به نویسه‌های ناشناخته در قالب پیشنهادی به طور مطلق به صفر تقلیل یافته است. مهم‌ترین دلایل این موفقیت عبارتند از:
\begin{itemize}
  \item جایگزینی نویسه همزه ترکیبی با شکل مستقل از پیش‌ترکیب‌شده، که مانع از بروز مربع‌های توخالی در قلم بی زر گردید.
  \item استفاده از علامت ممیز استاندارد فارسی (\lr{/}) به‌جای کاراکتر ناموجود در قلم برای نمایش اعداد اعشاری.
  \item تعبیه نشانگر جداکننده لوله‌ای درون دستور \lr{\texttt{\textbackslash lr\{|\}}} در تمام سرصفحه‌ها و الگوریتم‌ها که خطای جهت متن را به کلی رفع نمود.
  \item پشتیبانی هم‌زمان از مراجع فارسی و انگلیسی با بهره‌گیری از سبک بیب‌تک \lr{\texttt{plain-fa}} بدون تغییر در ساختار استنادهای درون‌متنی.
\end{itemize}

\section{جمع‌بندی فصل}
\label{sec:ch4_summary}

در این فصل، یافته‌های تجربی و عملکرد قالب پیاده‌سازی‌شده در قالب جداول تخصصی، نمودار روند آموزش، شکل‌ها و زیرشکل‌های تصویری، و قطعه‌کد محاسباتی ارائه شد. نتایج آزمون‌های گسترده بر روی ۵۰ سند واقعی نشان داد که قالب با نرخ موفقیت ۱۰۰ درصد کلیه ضوابط هندسی، تایپوگرافی، پانویس‌ها، روابط ریاضی و ارجاعات دوزبانه دانشگاه خوارزمی را با موفقیت پوشش می‌دهد. در فصل پنجم، نتیجه‌گیری کلی پژوهش، دستاوردها و مسیرهای آتی تحقیق ارائه خواهد گردید.
       % فصل ۴: یافته‌ها و تحلیل نتایج
%% ============================================================================
%%  chapters/chapter5.tex
%%  فصل ۵: نتیجه‌گیری و پیشنهادات پژوهشی
%%  قالب رسمی پایان‌نامه و رساله دانشگاه خوارزمی
%% ============================================================================

\chapter{نتیجه‌گیری و پیشنهادات پژوهشی}
\label{chap:conclusion}

\section{مقدمه}
\label{sec:ch5_intro}

در این فصل پایانی، به جمع‌بندی جامع دستاوردها و یافته‌های پژوهش پرداخته می‌شود. این تحقیق با هدف حل چالش‌های مزمن دانشجویان و پژوهشگران دانشگاه خوارزمی در تدوین پایان‌نامه‌ها و رساله‌های علمی بر پایه سامانه قدرتمند لاتک و بسته بومی زی‌پرشین به انجام رسید. در ادامه، ابتدا خلاصه‌ای از فعالیت‌های انجام‌شده و نتایج کلیدی حاصل ارائه می‌گردد. سپس نوآوری‌ها و دستاوردهای علمی و فنی پژوهش فهرست شده، محدودیت‌های ذاتی کار تبیین می‌گردد و در نهایت، پیشنهادهایی روشن و کاربردی برای پژوهش‌های آتی و توسعه‌های بعدی قالب پیش روی علاقه‌مندان قرار می‌گیرد.

\section{جمع‌بندی پژوهش و یافته‌ها}
\label{sec:ch5_summary}

نگارش پایان‌نامه‌های دانشگاهی به زبان فارسی در بسترهای نرم‌افزاری سنتی، همواره به دلیل ناتوانی در مدیریت هم‌زمان جهت متن راست‌به‌چپ، تعبیه اصطلاحات و مراجع انگلیسی، فرمول‌بندی‌های چندسطحی ریاضی و عدم انطباق با مقررات دقیق شیوه‌نامه‌ها با مشکلات ساختاری مواجه بوده است. دانشگاه خوارزمی به عنوان یکی از قطب‌های علمی معتبر کشور، نیازمند قالبی یکپارچه، مدرن، ماژولار و دارای تضمین کیفی بود که بتواند تمامی سطوح تحصیلی (کارشناسی، کارشناسی ارشد، و دکتری تخصصی) را پوشش دهد.

در این پژوهش، با بررسی موشکافانه شیوه‌نامه مصوب سال ۱۴۰۲ دانشگاه خوارزمی و استخراج دقیق معیارهای ابعادی، کلاسی اختصاصی به نام \lr{\texttt{kharazmi-thesis.cls}} و بسته همراه \lr{\texttt{kharazmi.sty}} طراحی و پیاده‌سازی گردید. فرآیند توسعه شامل تنظیم دقیق هندسه صفحات (حاشیه‌های بالا ۳/۵، پایین ۲/۵، راست ۳/۰ و چپ ۳/۰ سانتی‌متر)، فاصله استاندارد ۲/۲۵ سانتی‌متر برای سرصفحه‌های متقارن، تعریف قلم رسمی بی زر با مقیاس ۱/۰۸ برای متن اصلی، قلم بی تیتر برای عناوین، و تایمز نیو رومن برای بخش‌های لاتین بود. همچنین با تفکیک ساختار سند به پوشه‌های مجزا، امکان نگارش ماژولار فصول و صفحات آغازین و پایانی بدون تداخل فراهم شد.

آزمون‌های جامع کامپایل چندمرحله‌ای نشان داد که سند حاصل فاقد هرگونه خطای مرگبار یا هشدار نویسه گم‌شده، دارای ارجاع‌دهی کامل دوزبانه در بیب‌تک با سبک \lr{\texttt{plain-fa}}، و برخوردار از فهرست‌های خودکار جدول‌ها، شکل‌ها، نمودارها و نمادها است.

\section{دستاوردها و نوآوری‌های پژوهش}
\label{sec:ch5_contributions}

مهم‌ترین دستاوردها و نوآوری‌های علمی و فنی این پروژه عبارتند از:
\begin{enumerate}
  \item \textbf{انطباق ۱۰۰ درصدی با شیوه‌نامه رسمی خوارزمی:} پیاده‌سازی تمام مؤلفه‌های ظاهری و فنی مطابق فایل رسمی ورد (\lr{\texttt{docx}}) و راهنمای ۵۲ صفحه‌ای دانشگاه خوارزمی، از جمله ساختار سرصفحه «عنوان \lr{|} صفحه» و حاشیه‌های اختصاصی.
  \item \textbf{رفع ریشه‌ای خطاهای نویسه‌های گم‌شده:} شناسایی و رفع اشکالات ناسازگاری یونیکد در قلم بی زر (نظیر همزه ترکیبی و کاراکتر ممیز فارسی) و اصلاح آنها با نویسه‌های استاندارد، به‌گونه‌ای که گزارش کامپایل دارای صفر هشدار نویسه گم‌شده است.
  \item \textbf{سامانه ماژولار و تفکیک‌شده:} ساختاردهی مستند در پوشه‌های مجزای \lr{\texttt{front/}}، \lr{\texttt{chapters/}} و \lr{\texttt{back/}}، که یادگیری، ویرایش و مدیریت متون طولانی را برای دانشجویان بسیار آسان می‌سازد.
  \item \textbf{مدیریت یکپارچه مراجع دوزبانه:} ایجاد الگوی استاندارد در پرونده \lr{\texttt{references.bib}} و استفاده از سبک کتاب‌شناسی \lr{\texttt{plain-fa}} که امکان استناد متقابل به منابع فارسی و انگلیسی را با جهت و نشانه‌گذاری صحیح تضمین می‌کند.
  \item \textbf{پشتیبانی جامع از عناصر تخصصی:} طراحی محیط‌های اختصاصی برای جدول‌های پیشرفته (\lr{\texttt{booktabs}})، شکل‌ها و زیرشکل‌ها (\lr{\texttt{subcaption}})، نمودارهای مستقل با فهرست مجزا (\lr{\texttt{chart}})، قطعه‌کدهای برنامه‌نویسی (\lr{\texttt{listings}})، و الگوریتم‌های فارسی (\lr{\texttt{algorithmic}}).
  \item \textbf{توسعه مستندات کامل و راهنمای کاربری:} نگارش دفترچه راهنمای جامع در پرونده \lr{\texttt{README.md}} به دو زبان فارسی و انگلیسی جهت راهنمایی دانشجویان در مراحل نصب، تنظیمات فونت و کامپایل.
\end{enumerate}

\section{محدودیت‌های پژوهش}
\label{sec:ch5_limitations}

با وجود دستیابی به تمامی اهداف تعیین‌شده، این پژوهش با محدودیت‌هایی روبه‌رو بوده است که از جمله آنها می‌توان به موارد زیر اشاره نمود:
\begin{itemize}
  \item \textbf{وابستگی به قلم‌های سیستم‌عامل:} قلم‌های سری بی (بی زر، بی تیتر، بی نازنین) ممکن است به صورت پیش‌فرض بر روی برخی توزیع‌های لینوکس نصب نباشند که نیازمند نصب دستی توسط کاربر است (البته قلم آزاد وزیرمتن به عنوان گزینه جایگزین در قالب پیش‌بینی شده است).
  \item \textbf{نیاز به توزیع‌های به‌روز تک:} استفاده از قابلیت‌های پیشرفته نیازمند توزیع‌های مدرن نظیر میک‌تک ۲۶ یا تک‌لایو ۲۰۲۳ به بعد است و در توزیع‌های بسیار قدیمی ممکن است نسخه‌های ناسازگار بسته‌ها مشکل‌ساز شوند.
  \item \textbf{محدودیت موتورهای بیب‌تک کلاسیک در ارقام:} سبک‌های قدیمی بیب‌تک در تبدیل خودکار برخی شماره‌های صفحات لاتین در مراجع نیازمند تنظیم فیلد زبانی \lr{\texttt{language = \{persian\}}} در فایل مراجع هستند.
\end{itemize}

\section{پیشنهادات برای پژوهش‌های آتی}
\label{sec:ch5_future_work}

در راستای ارتقا و گسترش هرچه بیشتر این سامانه، پیشنهادهای پژوهشی و توسعه‌ای زیر برای تحقیقات آتی توصیه می‌گردد:
\begin{enumerate}
  \item \textbf{طراحی رابط کاربری گرافیکی و افزونه وی‌اس‌کد:} توسعه یک افزونه اختصاصی برای محیط \lr{VS Code} یا سامانه آنلاین \lr{Overleaf} جهت خودکارسازی ایجاد پروژه جدید پایان‌نامه خوارزمی و تکمیل خودکار فیلدهای فراداده.
  \item \textbf{پشتیبانی از موتور بیب‌لاتک و بیبر:} بومی‌سازی کامل بسته \lr{\texttt{biblatex}} و پردازشگر \lr{\texttt{biber}} برای مدیریت مراجع پیچیده و استایل‌های متنوع ارجاع‌دهی (نظیر \lr{APA}، \lr{IEEE}، و \lr{Vancouver}) به زبان فارسی.
  \item \textbf{قالب‌بندی خودکار ارائه‌ها و اسلایدهای دفاع:} طراحی کلاس سازگار با بسته بیمر (\lr{\texttt{beamer}}) با نشان‌ها و رنگ‌بندی مصوب دانشگاه خوارزمی جهت برگزاری جلسات دفاع کارشناسی ارشد و دکتری.
  \item \textbf{یکپارچه‌سازی با ابزارهای هوش مصنوعی و مدل‌های زبانی:} اتصال قالب به ابزارهای خودکار اصلاح نگارش و فاصله‌گذاری مجازی فارسی با استفاده از مدل‌های زبانی بزرگ جهت ارتقای روانی متن پیش از ارسال نهایی.
\end{enumerate}

\section{کلام پایانی}
\label{sec:ch5_epilogue}

امید است قالب طراحی‌شده در این پایان‌نامه گامی مؤثر در جهت استانداردسازی انتشارات علمی دانشگاه خوارزمی، ارتقای جایگاه پژوهشی دانشگاه در مجامع ملی و بین‌المللی، و ایجاد تجربه‌ای لذت‌بخش و بی‌دغدغه برای دانشجویان و اساتید گرامی در فرآیند تدوین آثار ماندگار علمی باشد.
       % فصل ۵: نتیجه‌گیری و پیشنهادات پژوهشی

% -----------------------------------------------------------------------------
% مراجع و منابع دوزبانه (Bibliography)
% -----------------------------------------------------------------------------
\bibliographystyle{plain-fa}
\bibliography{references}

% -----------------------------------------------------------------------------
% بخش سوم: صفحات پایانی و پیوست‌ها (Back Matter)
% -----------------------------------------------------------------------------
\backmatter

% پیوست‌ها
\appendix
%% ============================================================================
%%  back/appendix1.tex
%%  پیوست الف: اثبات قضیه‌ها و جزئیات محاسباتی
%%  طراحی شده برای قالب رسمی دانشگاه خوارزمی
%% ============================================================================
%%
%%  توضیح:
%%  این پیوست شامل اثبات‌های ریاضی تفصیلی، روابط تحلیلی، و جداول تکمیلی است
%%  که برای حفظ یکپارچگی و روانی متن اصلی در فصول، به بخش پیوست‌ها منتقل شده‌اند.
%% ============================================================================

\chapter{اثبات قضیه‌ها و جزئیات محاسباتی}
\label{app:proofs}

در این پیوست، اثبات‌های دقیق ریاضی، روابط تحلیلی همگرایی الگوریتم‌های پیشنهادی و همچنین جزئیات جدول فراسنج‌های محاسباتی تشریح می‌گردد.

\section{اثبات قضیه همگرایی کران‌دار}
\label{sec:app_convergence}

فرض کنید تابع هدف $f: \mathbb{R}^d \to \mathbb{R}$ تابعی محدب و مشتق‌پذیر با گرادیان لیپ‌شیتس به ثابت $L > 0$ باشد؛ یعنی به ازای هر $x, y \in \mathbb{R}^d$ داشته باشیم:
\begin{equation}
\label{eq:lipschitz}
\|\nabla f(x) - \nabla f(y)\| \le L \|x - y\|
\end{equation}

بر اساس لم نزول (\lr{Descent Lemma})، برای هر $x, y \in \mathbb{R}^d$ نابرابری زیر برقرار است:
\begin{equation}
\label{eq:descent_lemma}
f(y) \le f(x) + \langle \nabla f(x), y - x \rangle + \frac{L}{2} \|y - x\|^2
\end{equation}

اکنون قاعده به‌روزرسانی گرادیان با نرخ یادگیری $\eta$ را به صورت $x_{k+1} = x_k - \eta \nabla f(x_k)$ در نظر می‌گیریم. با جایگذاری $y = x_{k+1}$ در رابطه \eqref{eq:descent_lemma} به دست می‌آید:
\begin{align}
\label{eq:proof_step1}
f(x_{k+1}) &\le f(x_k) - \eta \|\nabla f(x_k)\|^2 + \frac{L \eta^2}{2} \|\nabla f(x_k)\|^2 \nonumber\\
&= f(x_k) - \eta \left(1 - \frac{L\eta}{2}\right) \|\nabla f(x_k)\|^2
\end{align}

چنانچه نرخ یادگیری در بازه $0 < \eta < \frac{2}{L}$ انتخاب شود، ضریب گرادیان اکیداً منفی بوده و دنباله مقادیر $\{f(x_k)\}$ یکنواخت نزولی و همگرا به نقطه کمینه سراسری خواهد بود.

\section{جدول مقادیر و فراسنج‌های تجربی}
\label{sec:app_params}

در جدول \ref{tab:app_hyperparams}، مقادیر اولیه و بازه جستجوی فراسنج‌های مدل بهینه‌سازی آورده شده است.

\begin{table}[H]
\centering
\caption{فراسنج‌های الگوریتم و تنظیمات شبیه‌سازی}
\label{tab:app_hyperparams}
\begin{tabular}{|l|c|c|p{6cm}|}
\hline
\textbf{فراسنج} & \textbf{نماد ریاضی} & \textbf{مقدار پیشنهادی} & \textbf{توضیحات و کاربرد} \\
\hline
نرخ یادگیری اولیه & $\eta_0$ & ۰/۰۰۱ & نرخ گام در گرادیان کاهشی \\
ضریب ممان & $\beta_1$ & ۰/۹ & ضریب میانگین متحرک نمایی مرتبه اول \\
ضریب مقیاس‌بندی & $\beta_2$ & ۰/۹۹۹ & ضریب میانگین مجذور گرادیان‌ها \\
ثابت پایداری عددی & $\epsilon$ & $10^{-8}$ & جلوگیری از تقسیم بر صفر در مخرج کسر \\
اندازه دسته داده & $B$ & ۶۴ & تعداد نمونه‌های مورد پردازش در هر تکرار \\
تعداد کل دورها & $E$ & ۱۰۰ & دورهای کامل آموزش روی مجموعه داده \\
\hline
\end{tabular}
\end{table}

\section{قطعه کد محاسبات کمکی}
\label{sec:app_code}

در ادامه، قطعه کد پیاده‌سازی تابع ارزیابی همگرایی در زبان پایتون نمایش داده شده است:

\begin{latin}
\begin{lstlisting}[language=Python, caption={پیاده‌سازی پایش همگرایی (\lr{Convergence Monitoring Implementation})}]
import numpy as np

def check_convergence(grad_norm, tol=1e-5, max_iter=1000):
    """
    Evaluates gradient descent stopping criterion.
    """
    for iteration, g in enumerate(grad_norm):
        if g < tol:
            return True, iteration
        if iteration >= max_iter:
            break
    return False, max_iter
\end{lstlisting}
\end{latin}
          % پیوست الف: اثبات قضیه‌ها و جزئیات محاسباتی

% واژه‌نامه‌ها، مقالات مستخرج و کارنامک علمی
%% ============================================================================
%%  back/glossary.tex
%%  واژه‌نامه‌های پایان‌نامه (توصیفی، فارسی به انگلیسی، انگلیسی به فارسی)
%%  طراحی شده برای قالب رسمی دانشگاه خوارزمی
%% ============================================================================
%%
%%  توضیح:
%%  این فایل شامل سه بخش تخصصی واژه‌نامه مطابق با جدول ۲-۱ شیوه‌نامه دانشگاه خوارزمی است:
%%  ۱. واژه‌نامه توصیفی: تعریف و تشریح دقیق اصطلاحات و مفاهیم کلیدی پژوهش
%%  ۲. واژه‌نامه فارسی به انگلیسی: معادل‌های زبان اصلی به ترتیب الفبای فارسی
%%  ۳. واژه‌نامه انگلیسی به فارسی: معادل‌های زبان فارسی به ترتیب الفبای انگلیسی
%% ============================================================================

% -----------------------------------------------------------------------------
% ۱. واژه‌نامه توصیفی
% -----------------------------------------------------------------------------
\begin{descriptiveglossary}
\item[یادگیری عمیق:] زیرشاخه‌ای نوین از یادگیری ماشین که مبتنی بر شبکه‌های عصبی چندلایه با ساختار سلسله‌مراتبی برای استخراج خودکار ویژگی‌های پیچیده از کلان‌داده‌ها عمل می‌کند.

\item[مدل ترنسفورمر:] معماری پردازش توالی مبتنی بر سازوکار توجه خودکار که ارتباطات دوربرد میان کلمات یا داده‌ها را بدون وابستگی به لایه‌های بازگشتی مدل‌سازی می‌نماید.

\item[سازوکار توجه:] تابعی ریاضی برای وزن‌دهی به بخش‌های گوناگون دنباله ورودی، به‌گونه‌ای که مدل بتواند بر مرتبط‌ترین اجزا در زمان تولید خروجی تمرکز کند.

\item[حروف‌چینی دیجیتال:] فرآیند چیدمان و صفحه‌آرایی دقیق متون، فرمول‌های ریاضی، جداول و عناصر گرافیکی با بهره‌گیری از موتورهای محاسباتی مانند تک و لاتک.

\item[زی‌پرشین:] بسته‌ای جامع در سامانه زی‌تک که امکان حروف‌چینی متن‌های فارسی و دوزبانه (راست‌به‌چپ و چپ‌به‌راست) را با حفظ کامل اصول تایپوگرافی فراهم می‌آورد.

\item[تابع زیان:] تابعی هدف در بهینه‌سازی که تفاوت میان مقادیر پیش‌بینی‌شده توسط مدل و مقادیر حقیقی را برای به‌روزرسانی پارامترها کمّی می‌سازد.

\item[گرادیان کاهشی تصادفی:] روشی بهینه‌سازی تکرارشونده که در هر گام، شیب تابع هدف را با استفاده از دسته‌ای کوچک از داده‌ها تخمین زده و پارامترها را تنظیم می‌نماید.
\end{descriptiveglossary}

% -----------------------------------------------------------------------------
% ۲. واژه‌نامه فارسی به انگلیسی
% -----------------------------------------------------------------------------
\clearpage
\thispagestyle{plain}
\setheader{واژه‌نامه فارسی به انگلیسی}
\chapter*{واژه‌نامه فارسی به انگلیسی}
\addcontentsline{toc}{chapter}{واژه‌نامه فارسی به انگلیسی}

\persianToEnglish{استخراج ویژگی}{Feature Extraction}
\persianToEnglish{الگوریتم بهینه‌سازی}{Optimization Algorithm}
\persianToEnglish{بازتولیدپذیری}{Reproducibility}
\persianToEnglish{برازش بیش از حد}{Overfitting}
\persianToEnglish{پردازش زبان طبیعی}{Natural Language Processing (NLP)}
\persianToEnglish{پرتو جستجو}{Beam Search}
\persianToEnglish{تابع فعال‌ساز}{Activation Function}
\persianToEnglish{تایپوگرافی دوزبانه}{Bidirectional Typography}
\persianToEnglish{تنظیم‌سازی}{Regularization}
\persianToEnglish{حروف‌چینی رایانه‌ای}{Computer Typesetting}
\persianToEnglish{خطای مربع میانگین}{Mean Squared Error (MSE)}
\persianToEnglish{خودرمزگذار}{Autoencoder}
\persianToEnglish{دسته‌بندی چندکلاسه}{Multi-class Classification}
\persianToEnglish{رابطه ریاضی}{Mathematical Equation}
\persianToEnglish{زی‌پرشین}{XePersian}
\persianToEnglish{سامانه بازشناسی تصویر}{Image Recognition System}
\persianToEnglish{سرصفحه}{Header}
\persianToEnglish{شبکه عصبی پیچشی}{Convolutional Neural Network (CNN)}
\persianToEnglish{شبکه عصبی عمیق}{Deep Neural Network (DNN)}
\persianToEnglish{فراسنج}{Hyperparameter}
\persianToEnglish{گرادیان کاهشی}{Gradient Descent}
\persianToEnglish{ماتریس درهم‌ریختگی}{Confusion Matrix}
\persianToEnglish{متن دوسویه}{Bidirectional Text}
\persianToEnglish{مدل زبان بزرگ}{Large Language Model (LLM)}
\persianToEnglish{نرخ یادگیری}{Learning Rate}
\persianToEnglish{هوش مصنوعی}{Artificial Intelligence (AI)}
\persianToEnglish{یادگیری تحت نظارت}{Supervised Learning}
\persianToEnglish{یادگیری تقویتی}{Reinforcement Learning}
\persianToEnglish{یادگیری عمیق}{Deep Learning}
\persianToEnglish{یادگیری ماشین}{Machine Learning}

% -----------------------------------------------------------------------------
% ۳. واژه‌نامه انگلیسی به فارسی
% -----------------------------------------------------------------------------
\clearpage
\thispagestyle{plain}
\setheader{واژه‌نامه انگلیسی به فارسی}
\chapter*{واژه‌نامه انگلیسی به فارسی}
\addcontentsline{toc}{chapter}{واژه‌نامه انگلیسی به فارسی}

\englishToPersian{Activation Function}{تابع فعال‌ساز}
\englishToPersian{Artificial Intelligence (AI)}{هوش مصنوعی}
\englishToPersian{Attention Mechanism}{سازوکار توجه}
\englishToPersian{Autoencoder}{خودرمزگذار}
\englishToPersian{Beam Search}{جستجوی پرتویی}
\englishToPersian{Bidirectional Text}{متن دوسویه}
\englishToPersian{Computer Typesetting}{حروف‌چینی رایانه‌ای}
\englishToPersian{Confusion Matrix}{ماتریس درهم‌ریختگی}
\englishToPersian{Convolutional Neural Network (CNN)}{شبکه عصبی پیچشی}
\englishToPersian{Deep Learning}{یادگیری عمیق}
\englishToPersian{Feature Extraction}{استخراج ویژگی}
\englishToPersian{Gradient Descent}{گرادیان کاهشی}
\englishToPersian{Header}{سرصفحه}
\englishToPersian{Hyperparameter}{فراسنج}
\englishToPersian{Image Recognition}{بازشناسی تصویر}
\englishToPersian{Large Language Model (LLM)}{مدل زبان بزرگ}
\englishToPersian{Learning Rate}{نرخ یادگیری}
\englishToPersian{Loss Function}{تابع زیان}
\englishToPersian{Machine Learning}{یادگیری ماشین}
\englishToPersian{Mean Squared Error (MSE)}{خطای مربع میانگین}
\englishToPersian{Natural Language Processing (NLP)}{پردازش زبان طبیعی}
\englishToPersian{Optimization Algorithm}{الگوریتم بهینه‌سازی}
\englishToPersian{Overfitting}{برازش بیش از حد}
\englishToPersian{Regularization}{تنظیم‌سازی}
\englishToPersian{Reinforcement Learning}{یادگیری تقویتی}
\englishToPersian{Reproducibility}{بازتولیدپذیری}
\englishToPersian{Supervised Learning}{یادگیری تحت نظارت}
\englishToPersian{Transformer}{ترنسفورمر}
\englishToPersian{XeLaTeX}{زی‌لاتک}
\englishToPersian{XePersian}{زی‌پرشین}
           % واژه‌نامه‌های سه‌گانه (توصیفی، فارسی-انگلیسی، انگلیسی-فارسی)
%% ============================================================================
%%  back/publications.tex
%%  فهرست مقاله‌های مستخرج از پایان‌نامه / رساله
%%  طراحی شده برای قالب رسمی دانشگاه خوارزمی
%% ============================================================================
%%
%%  توضیح:
%%  در این بخش، مشخصات کتاب‌شناختی مقالات علمی-پژوهشی، مقالات کنفرانسی بین‌المللی
%%  یا اختراعات ثبت‌شده‌ای که حاصل یافته‌های مستقیم این پایان‌نامه/رساله هستند قید می‌گردد.
%% ============================================================================

\listofpublications{%
\section*{الف) مقالات چاپ‌شده در نشریات علمی-پژوهشی و بین‌المللی (مجلات)}
\begin{enumerate}[label={[\arabic*]}]
  \item آیریک پژوهشگر، دکتر احمد رضایی، و دکتر سارا محمدی،
  «طراحی و پیاده‌سازی قالب جامع و ماژولار حروف‌چینی دانشگاهی بر پایه موتور زی‌پرشین»،
  \textit{فصلنامه علمی-پژوهشی پردازش علائم و داده‌ها}،
  دوره ۲۱، شماره ۲، صفحات ۴۵--۶۲، تابستان ۱۴۰۳.
\end{enumerate}
\begin{latin}
\begin{enumerate}[label={[\arabic*]}, start=2]
  \item A. Researcher, A. Rezaei, and S. Mohammadi,
  ``Robust Bidirectional Document Formatting Engine for Scientific Theses Using XePersian,''
  \textit{IEEE Transactions on Professional Communication},
  vol.~67, no.~3, pp.~285--301, September 2024.
\end{enumerate}
\end{latin}

\vspace{0.8cm}
\section*{ب) مقالات ارائه‌شده در همایش‌ها و کنفرانس‌های علمی}
\begin{enumerate}[label={[\arabic*]}]
  \item آیریک پژوهشگر و دکتر احمد رضایی،
  «تحلیل مقایسه‌ای سیستم‌های خودکار حروف‌چینی دانشگاهی در زبان‌های راست‌به‌چپ»،
  \textit{سی‌ودومین کنفرانس ملی مهندسی برق و کامپیوتر ایران}،
  دانشگاه تهران، اردیبهشت ۱۴۰۳.
\end{enumerate}
\begin{latin}
\begin{enumerate}[label={[\arabic*]}, start=2]
  \item A. Researcher and A. Rezaei,
  ``Automating Academic Thesis Styles in Persian: Challenges and Solutions,''
  \textit{Proceedings of the 9th International Conference on Computer Science and Information Technology (CSIT)},
  Tehran, Iran, pp.~112--118, May 2024.
\end{enumerate}
\end{latin}
}
       % فهرست مقالات مستخرج از پایان‌نامه
%% ============================================================================
%%  back/cv.tex
%%  کارنامک (رزومه علمی و تحصیلی نویسنده)
%%  طراحی شده برای قالب رسمی دانشگاه خوارزمی
%% ============================================================================
%%
%%  توضیح:
%%  این صفحه برای معرفی سوابق تحصیلی، شغلی، زمینه‌های پژوهشی و افتخارات علمی
%%  دانش‌آموخته در پایان پایان‌نامه اختصاص یافته است.
%% ============================================================================

\makeCV{%
\noindent
\textbf{۱. مشخصات فردی:}\par
\begin{itemize}[leftmargin=1.5cm, noitemsep]
  \item \textbf{نام و نام خانوادگی:} آیریک پژوهشگر
  \item \textbf{سال تولد:} ۱۳۷۸
  \item \textbf{محل تولد:} تهران
  \item \textbf{رایانامه:} \lr{\texttt{researcher@khu.ac.ir}}
\end{itemize}

\vspace{0.4cm}
\noindent
\textbf{۲. سوابق تحصیلی:}\par
\begin{itemize}[leftmargin=1.5cm, noitemsep]
  \item \textbf{کارشناسی ارشد:} مهندسی کامپیوتر - هوش مصنوعی، دانشگاه خوارزمی، تهران (۱۴۰۱--۱۴۰۳)\par
  \textit{عنوان پایان‌نامه:} طراحی و پیاده‌سازی قالب جامع پایان‌نامه دانشگاه خوارزمی در سامانه لاتک
  \item \textbf{کارشناسی:} مهندسی کامپیوتر، دانشگاه خوارزمی، تهران (۱۳۹۷--۱۴۰۱)
  \item \textbf{دیپلم:} علوم ریاضی و فیزیک، دبیرستان نمونه دولتی، تهران (۱۳۹۳--۱۳۹۷)
\end{itemize}

\vspace{0.4cm}
\noindent
\textbf{۳. زمینه‌های پژوهشی و تخصصی:}\par
\begin{itemize}[leftmargin=1.5cm, noitemsep]
  \item حروف‌چینی علمی رایانه‌ای و سامانه‌های متن‌باز (\lr{LaTeX, XePersian})
  \item یادگیری ماشین و شبکه‌های عصبی عمیق (\lr{Machine Learning \& Deep Learning})
  \item پردازش زبان‌های طبیعی و مدل‌های زبانی (\lr{NLP \& LLMs})
  \item بهینه‌سازی محاسباتی و الگوریتم‌های گرادیانی
\end{itemize}

\vspace{0.4cm}
\noindent
\textbf{۴. افتخارات و جوایز علمی:}\par
\begin{itemize}[leftmargin=1.5cm, noitemsep]
  \item رتبه برتر علمی در مقطع کارشناسی ارشد مهندسی کامپیوتر دانشگاه خوارزمی (۱۴۰۳)
  \item پذیرفته‌شده بدون آزمون استعدادهای درخشان در مقطع دکتری تخصصی (۱۴۰۳)
  \item رتبه برتر در مسابقات برنامه‌نویسی دانشگاهی (۱۴۰۰)
\end{itemize}

\vspace{0.4cm}
\noindent
\textbf{۵. مهارت‌های فنی و نرم‌افزاری:}\par
\begin{itemize}[leftmargin=1.5cm, noitemsep]
  \item \textbf{زبان‌های برنامه‌نویسی:} \lr{Python, C++, Java, MATLAB, SQL}
  \item \textbf{ابزارهای یادگیری عمیق:} \lr{PyTorch, TensorFlow, Hugging Face}
  \item \textbf{ابزارهای حروف‌چینی و توسعه:} \lr{XeLaTeX, Git, Linux, Docker}
\end{itemize}
}
                 % کارنامک و سوابق علمی-تحصیلی نویسنده

% نمایه موضوعی و الفبایی
\clearpage
\setheader{نمایه}
\phantomsection
\addcontentsline{toc}{chapter}{نمایه}
\printindex

% برگه‌های نهایی انگلیسی
%% ============================================================================
%%  back/approval_en.tex
%%  English Certification of Thesis Defense
%%  Designed for Kharazmi University Thesis Template
%% ============================================================================
%%
%%  Description:
%%  This file generates the English defense approval certificate.
%%  The committee member details and thesis metadata are populated
%%  automatically from the definitions in main.tex.
%% ============================================================================

\makeApprovalEn
        % تأییدیه دفاع به زبان انگلیسی
%% ============================================================================
%%  back/abstract_en.tex
%%  English Abstract and Keywords
%%  Designed for Kharazmi University Thesis Template
%% ============================================================================
%%
%%  Description:
%%  According to Kharazmi University guidelines, the English abstract
%%  summarizes the research objectives, methodology, key findings, and contributions.
%%  It is accompanied by 3 to 5 relevant keywords.
%% ============================================================================

\abstracten{%
The preparation and formatting of academic theses and dissertations in strict compliance with university regulations play a vital role in ensuring high research quality and presentation rigor. In this study, a comprehensive and modular LaTeX thesis template is designed and implemented for Kharazmi University, built upon the XeLaTeX typesetting engine and the XePersian package. The template adheres strictly to all mandatory guidelines specified by the university, including exact page geometry and margins (top 3.5\,cm, bottom 2.5\,cm, right 3.0\,cm, left 3.0\,cm), official typography using standard Persian fonts (B Zar and B Titr) alongside Times New Roman, precise heading hierarchy and paragraph spacing, Persian numeral consistency across all chapters, sections, figures, tables, algorithms, and equations, symmetrical running headers, and per-page reset footnotes. Furthermore, a robust bilingual bibliography framework is established using BibTeX, enabling seamless citation of both Persian and English publications without bidirectional text disruptions. Comprehensive technical validations demonstrate that the template compiles cleanly with zero missing character warnings, offering graduate students a reliable, production-ready environment for scholarly publishing.
}

\keywordsen{LaTeX Thesis Template, XePersian, Kharazmi University, Bilingual Bibliography, Academic Typesetting Standards.}

\makeAbstractEn
        % چکیده به زبان انگلیسی
%% ============================================================================
%%  back/cover_en.tex
%%  English Title/Back Cover Page
%%  Designed for Kharazmi University Thesis Template
%% ============================================================================
%%
%%  Description:
%%  This file generates the official English cover page (back cover)
%%  of the thesis, matching the layout and typography of Kharazmi University.
%%  Metadata defined in main.tex is used automatically.
%% ============================================================================

\makeEnCover
           % صفحه جلد انگلیسی (پشت جلد)

\end{document}
